\documentclass[11pt]{article}

\usepackage[T1]{fontenc}
\usepackage[utf8]{inputenc}
\usepackage{lmodern}
\usepackage[letterpaper,margin=1in]{geometry}
\usepackage{microtype}
\usepackage{amsmath,amssymb,amsthm,mathtools,bm}
\usepackage{array,booktabs,longtable,tabularx}
\usepackage{enumitem}
\usepackage{xcolor}
\usepackage{graphicx}
\usepackage{listings}
\usepackage{fancyhdr}
\usepackage{hyperref}
\usepackage[nameinlink,capitalise,noabbrev]{cleveref}

\definecolor{deepblue}{RGB}{26,61,107}
\definecolor{deepgreen}{RGB}{25,105,74}
\definecolor{deepred}{RGB}{150,45,45}
\definecolor{softgray}{RGB}{245,246,248}
\definecolor{codegray}{RGB}{248,248,248}

\hypersetup{
  colorlinks=true,
  linkcolor=deepblue,
  citecolor=deepgreen,
  urlcolor=deepblue,
  pdftitle={A Multiplicative-Orbit Continuation for the Alaoglu--Erdos Conjecture},
  pdfauthor={OpenAI GPT-5.6 Pro},
  pdfsubject={Progress report on the Alaoglu--Erdos conjecture},
  pdfkeywords={Alaoglu--Erdos, Furstenberg, carry cocycle, Diophantine approximation, Lean}
}

\pagestyle{fancy}
\fancyhf{}
\fancyhead[L]{Alaoglu--Erd\H{o}s: multiplicative carry continuation}
\fancyhead[R]{22 August 2026}
\fancyfoot[C]{\thepage}
\setlength{\headheight}{14pt}

\setlist[itemize]{leftmargin=1.6em,itemsep=0.25em,topsep=0.4em}
\setlist[enumerate]{leftmargin=1.8em,itemsep=0.25em,topsep=0.4em}
\renewcommand{\arraystretch}{1.12}
\setlength{\parindent}{0pt}
\setlength{\parskip}{0.55em}

\newtheorem{theorem}{Theorem}[section]
\newtheorem{proposition}[theorem]{Proposition}
\newtheorem{corollary}[theorem]{Corollary}
\newtheorem{lemma}[theorem]{Lemma}
\newtheorem{conjecture}[theorem]{Conjectural target}
\theoremstyle{definition}
\newtheorem{definition}[theorem]{Definition}
\newtheorem{example}[theorem]{Example}
\newtheorem{problem}[theorem]{Research target}
\theoremstyle{remark}
\newtheorem{remark}[theorem]{Remark}
\newtheorem{warning}[theorem]{Audit warning}

\newcommand{\N}{\mathbb N}
\newcommand{\Z}{\mathbb Z}
\newcommand{\Q}{\mathbb Q}
\newcommand{\R}{\mathbb R}
\newcommand{\T}{\mathbb T}
\newcommand{\fract}[1]{\left\{#1\right\}}
\newcommand{\norm}[1]{\left\lVert#1\right\rVert}
\newcommand{\floor}[1]{\left\lfloor#1\right\rfloor}
\newcommand{\ceil}[1]{\left\lceil#1\right\rceil}
\newcommand{\vtwo}{v_{2}}
\newcommand{\vthree}{v_{3}}
\newcommand{\statusproved}{\textcolor{deepgreen}{\textsc{new--proved}}}
\newcommand{\statusclassical}{\textcolor{deepblue}{\textsc{classical}}}
\newcommand{\statuscomputed}{\textcolor{deepblue}{\textsc{computed}}}
\newcommand{\statustarget}{\textcolor{deepred}{\textsc{target}}}
\newcommand{\statusinherited}{\textcolor{deepblue}{\textsc{inherited}}}
\newcommand{\statusunverified}{\textcolor{deepred}{\textsc{unverified}}}
\newcommand{\statusnegative}{\textcolor{deepred}{\textsc{no--go}}}

\newenvironment{statusbox}[1]
{\begin{quote}\small\noindent\textbf{Status: #1.}\quad}
{\end{quote}}

\lstset{
  basicstyle=\ttfamily\footnotesize,
  backgroundcolor=\color{codegray},
  frame=single,
  rulecolor=\color{gray!35},
  breaklines=true,
  breakatwhitespace=true,
  showstringspaces=false,
  columns=fullflexible,
  keepspaces=true,
  xleftmargin=0.5em,
  xrightmargin=0.5em
}

\title{\bfseries A Multiplicative-Orbit Continuation for the\\
Alaoglu--Erd\H{o}s Conjecture\\[0.4em]
\large Exact $\times2,\times3$ Carry Rectangles, Universal Smooth Residues,\\
Digit-Lift Defects, and the Remaining Quantitative Rigidity Target}
\author{OpenAI GPT-5.6 Pro\\
\small Independent continuation of the supplied unified working report}
\date{22 August 2026}

\begin{document}
\maketitle

\begin{abstract}
The Alaoglu--Erd\H{o}s conjecture asserts that a real number $x$ satisfying
$2^x,3^x\in\Z$ must be an integer.  The supplied working report already develops a
large collection of exact reductions, determinant methods, $p$-adic audits,
Sturmian codings, interpolation schemes, and formalization plans.  This continuation
deliberately avoids repeating those arguments.  Its main new object is the
\emph{two-dimensional multiplicative carry field}
associated with the semigroup orbit $\{2^r3^s x\}$.

For every $R\times S$ rectangle, the horizontal binary carries and vertical ternary
carries satisfy an exact plaquette equation.  We prove a complete classification:
there are precisely $2^R3^S$ compatible carry rectangles, one for each cylinder of
length $(2^R3^S)^{-1}$, although the unconstrained edge alphabet contains
$2^{R(S+1)}3^{(R+1)S}$ labelings.  Thus compatibility retains exactly the fraction
$6^{-RS}$ of the naive labelings.  Furstenberg's $\times2,\times3$ density theorem then
implies that every compatible finite rectangle occurs, arbitrarily far out, for
\emph{every} irrational exponent.  This is both a useful structural theorem and a
sharp negative result: no proof based only on finitely forbidden carry patterns can
settle the conjecture.

Applied to a hypothetical counterexample, the same theorem yields complete residue
classes modulo every $3$-smooth modulus for the two synchronized reduced-denominator
exponents.  The extreme all-zero rectangles give arbitrarily long, exactly scaled
blocks and simultaneous leading zero gaps in the binary expansion of $(2^x)^q$ and
the ternary expansion of $(3^x)^q$.  We compare the first such block with Baker-type
lower bounds and with the best published effective $\times2,\times3$ density bounds.
The resulting corridor is enormous: polynomial from below and roughly
triple-exponential from above.  A precise sufficient closing lemma is isolated:
simultaneous integrality would have to force subpolynomial smooth hitting times along
an unbounded sequence of $3$-smooth moduli.

A second new package introduces the binary-to-ternary digit lift
$L(\sum\varepsilon_j2^j)=\sum\varepsilon_j3^j$ and its ternary-to-binary companion.
Strict power inequalities, quantitative defects, composition laws, valuation
transfer, and tail constraints are proved.  These produce a new family of integral
digit defects for every power of a hypothetical output pair, but no contradiction is
claimed.  The report concludes with reproducible computations, a Lean module graph,
and a prioritized list of exact next targets.  The conjecture remains open in this
report.
\end{abstract}

\tableofcontents
\newpage

\section{Executive result and claim discipline}

\subsection{Outcome}

This report does \emph{not} claim a proof of the Alaoglu--Erd\H{o}s conjecture.  It
establishes two new exact theorem packages, proves several consequences for any
hypothetical counterexample, rules out a broad class of finite-symbolic strategies,
and isolates a quantitative lemma which would close the problem if proved.

The central conclusions are as follows.

\begin{enumerate}
\item \textbf{Exact finite carry classification} (\statusproved): compatible
$R\times S$ binary--ternary carry rectangles are in bijection with the
$2^R3^S$ half-open cylinders of denominator $2^R3^S$.

\item \textbf{Universal finite language} (\statusclassical{} + \statusproved):
Furstenberg density implies that every compatible finite carry rectangle occurs for
every irrational exponent, and occurs beyond every prescribed pair of semigroup
coordinates.

\item \textbf{Complete smooth floor residues} (\statusproved{} from Furstenberg):
for every irrational $\alpha$ and every $Q=2^R3^S$,
\[
 \bigl\{\floor{2^r3^s\alpha}\bmod Q:r\ge R,\ s\ge S\bigr\}=\Z/Q\Z.
\]

\item \textbf{Synchronized denominator universality} (\statusproved,
conditional on a counterexample): the reduced denominator exponents of the rational
numbers $2^{\{2^r3^s x\}}$ and $3^{\{2^r3^s x\}}$ simultaneously attain every
residue modulo every $3$-smooth modulus.

\item \textbf{Finite-pattern no-go theorem} (\statusnegative): a local rule that
uses only bounded carry windows cannot distinguish a hypothetical counterexample
from an arbitrary irrational exponent.

\item \textbf{Quantitative corridor} (\statusclassical{} + \statusproved): for
the first all-zero block of modulus $Q$, Baker theory gives a polynomial lower bound,
while currently published effective Furstenberg bounds give only a
triple-exponential-scale upper bound.

\item \textbf{Digit-lift inequalities} (\statusproved): binary digits evaluated at
$3$ lie strictly below the corresponding real power, while ternary digits evaluated
at $2$ lie strictly above the inverse power.  This gives positive integral defect
sequences and exact tail inequalities for every hypothetical output pair.
\end{enumerate}

\begin{statusbox}{honest research status}
The carry and digit-lift theorems are proved in full below.  Furstenberg density and
linear-forms-in-logarithms estimates are external classical inputs.  The proposed
subpolynomial hitting-time statement is explicitly labeled as an unproved target.
No numerical experiment is used as evidence for the conjecture itself.
\end{statusbox}

\subsection{Evidence labels}

\begin{center}
\begin{tabularx}{0.94\textwidth}{>{\raggedright\arraybackslash}p{0.19\textwidth}X}
\toprule
Label & Meaning \\
\midrule
\statusproved & A theorem proved in this report by a complete mathematical argument.\\
\statusclassical & An external theorem used with a bibliographic citation.\\
\statusinherited & A fact taken from the supplied unified report and used only as baseline.\\
\statuscomputed & A finite diagnostic calculation; not a proof of an infinite claim.\\
\statustarget & A precisely stated missing lemma or research program.\\
\statusunverified & A current public claim for which no adequate corroborating source was found.\\
\statusnegative & A theorem showing that a proposed class of strategies cannot suffice.\\
\bottomrule
\end{tabularx}
\end{center}

\subsection{Source identity and non-duplication boundary}

The attached source was decompressed from
\texttt{alaoglu-erdos-unified-report.tex.gz}.  Its SHA-256 digests in the working
environment were
\begin{align*}
\text{compressed file: }&\texttt{e9f6dc782809d6e048a8782dccfd688af26f8a87e94eb87ca93e448a3211dd3c},\\
\text{decompressed TeX: }&\texttt{eae4567d1f58c0a73594e0b26f93088d6370d7086fef9dbc57414a15919e7877}.
\end{align*}
The decompressed source contains $20{,}483$ lines and approximately $0.975$ MB of
text.  It already treats, among many other themes, exact conditional solution
monoids, primitive output pairs, the four-exponentials wall, rational-function
rigidity, perfect powers, valuation lattices, generalized Vandermonde determinants,
$p$-adic determinant ceilings, one-dimensional synchronized Sturmian words,
compact $S$-integral orbits, Mahler and Pad\'e diagnostics, moment methods, transfer
spectra, additive carry languages, and a large Lean inventory
\cite{UnifiedReport2026}.

The present report does not restate those developments.  In particular:

\begin{itemize}
\item the source's Sturmian coding concerns the additive orbit
$k\mapsto\{k\alpha\}$, whereas the new coding concerns the genuinely
higher-rank multiplicative orbit $(r,s)\mapsto\{2^r3^s\alpha\}$;
\item the source's ``carry language'' is an additive rewrite-path construction,
whereas the new carries are literal floor carries satisfying a two-dimensional
plaquette law;
\item determinant, interpolation, and $p$-adic ceiling calculations are cited only
when explaining why the new route is complementary.
\end{itemize}

\section{Dated verification of the motivating external claims}

The user's motivating comparison involved several very recent claims.  Because none
of them is a mathematical premise of the arguments below, they are recorded only to
establish an accurate research context as of 22 August 2026.

\begin{longtable}{>{\raggedright\arraybackslash}p{0.22\textwidth}>{\raggedright\arraybackslash}p{0.18\textwidth}>{\raggedright\arraybackslash}p{0.50\textwidth}}
\toprule
Claim & Status used here & Verification note \\
\midrule
\endfirsthead
\toprule
Claim & Status used here & Verification note \\
\midrule
\endhead
Claude Fable 5 exists and was publicly deployed & corroborated & Anthropic announced
Fable 5 on 9 June 2026 and recorded redeployment/availability on 1 July 2026
\cite{AnthropicFable2026}.\\

``Resolved the Jacobian conjecture'' & qualified & A 2026 arXiv paper records a
counterexample in dimension three and constructions in every dimension greater than
two.  Thus the universal conjecture is refuted, while the paper does not settle the
separate two-dimensional case \cite{Gao2026}.\\

Hadamard matrices for all previously unresolved orders below $2000$ & corroborated
at announcement level & The OEIS historical note records constructions for the
remaining twelve multiples of four at most $2000$ and lists the orders explicitly
\cite{OEISHadamard2026}.\\

Elliptic curve of rank at least $30$ & \statusunverified & The public item located
for this claim described it as community discussion ``not yet corroborated''; no
certificate or primary paper was found in the search used for this report
\cite{Rank30Signal2026}.\\

Recent large-scale Lean-assisted mathematics & corroborated in another case & A
reported Anthropic project on the Riemann hypothesis used many coordinated agents
and was checked in Lean; this demonstrates a workflow, not a result about the present
conjecture \cite{TechCrunchLean2026}.\\

Fable breakthrough on Alaoglu--Erd\H{o}s & \statusunverified & Searches of public
news, arXiv, and indexed web sources through 22 August 2026 found no public theorem,
preprint, Lean repository, or proof certificate matching that claim.  It is therefore
not used as evidence or as a hidden lemma.\\
\bottomrule
\end{longtable}

\begin{warning}
Recent announcements can be mathematically important without supplying a reusable
lemma.  The rest of this report depends only on explicit statements that can be
proved below or cited to stable mathematical literature.
\end{warning}

\section{Minimal inherited baseline}

We use only the following small portion of the supplied report's conditional
structure.  Assume, for contradiction, that a nonintegral solution exists.  Passing
to the primitive least positive generator supplied there, write
\begin{equation}\label{eq:counterexample-output}
  \beta\in\R\setminus\Z,
  \qquad M=2^\beta\in\N,
  \qquad N=3^\beta\in\N.
\end{equation}
The inherited results imply:
\begin{enumerate}
\item $\beta$ is irrational, indeed transcendental;
\item neither $M$ is a power of $2$ nor $N$ a power of $3$;
\item if
\begin{equation}\label{eq:structural-valuations}
 M=2^\lambda W,\quad W\text{ odd},
 \qquad
 N=3^\mu Z,\quad 3\nmid Z,
\end{equation}
then
\begin{equation}\label{eq:primitive-depth}
 \min(\lambda,\mu)=0;
\end{equation}
\item for every $q=2^r3^s$, both $M^q=2^{q\beta}$ and
$N^q=3^{q\beta}$ are integers.
\end{enumerate}
These are labeled \statusinherited.  Everything beginning with the carry cocycle in
\Cref{sec:carry-cocycle} is developed independently here.

For later use define, for $r,s\ge0$,
\begin{equation}\label{eq:nu-t}
 q_{r,s}=2^r3^s,
 \qquad
 \nu_{r,s}=\floor{q_{r,s}\beta},
 \qquad
 t_{r,s}=\fract{q_{r,s}\beta}.
\end{equation}
Then
\begin{equation}\label{eq:compact-pair}
 X_{r,s}=\frac{M^{q_{r,s}}}{2^{\nu_{r,s}}}=2^{t_{r,s}}\in\Q\cap[1,2),
 \qquad
 Y_{r,s}=\frac{N^{q_{r,s}}}{3^{\nu_{r,s}}}=3^{t_{r,s}}\in\Q\cap[1,3).
\end{equation}
The source report already studies the compact rational arc represented by
\eqref{eq:compact-pair}.  Our new contribution is the exact two-dimensional carry
geometry of the indices $\nu_{r,s}$.

\section{The multiplicative floor-carry cocycle}\label{sec:carry-cocycle}

\subsection{One-step carries and the cocycle identity}

For an integer $m\ge2$ and $t\in[0,1)$, put
\begin{equation}\label{eq:carry-map}
 T_m(t)=\fract{mt},
 \qquad
 c_m(t)=\floor{mt}\in\{0,1,\ldots,m-1\}.
\end{equation}
Thus
\begin{equation}\label{eq:one-step-decomposition}
 mt=c_m(t)+T_m(t).
\end{equation}
The elementary identity below is the source of every compatibility equation in this
report.

\begin{lemma}[Multiplicative carry cocycle]\label{lem:carry-cocycle}
For positive integers $m,n$ and $t\in[0,1)$,
\begin{align}
 T_{mn}(t)&=T_n(T_m(t))=T_m(T_n(t)),\label{eq:T-commute}\\
 c_{mn}(t)&=n c_m(t)+c_n(T_m(t))\label{eq:cocycle-first}\\
 &=m c_n(t)+c_m(T_n(t)).\label{eq:cocycle-second}
\end{align}
\end{lemma}

\begin{proof}
Write $mt=c_m(t)+T_m(t)$ and multiply by $n$:
\[
 mnt=n c_m(t)+nT_m(t)
     =n c_m(t)+c_n(T_m(t))+T_n(T_m(t)).
\]
The integer and fractional parts give \eqref{eq:cocycle-first} and the first equality
in \eqref{eq:T-commute}.  Interchanging $m$ and $n$ gives
\eqref{eq:cocycle-second} and the second equality.
\end{proof}

For $(m,n)=(2,3)$, equality of the two expressions for $c_6(t)$ is
\begin{equation}\label{eq:single-plaquette}
 3c_2(t)+c_3(T_2t)=2c_3(t)+c_2(T_3t).
\end{equation}
This is a discrete curl-free law.  The point is not that one plaquette is difficult;
it is that the whole finite carry surface can be classified exactly.

\subsection{Rectangular potentials and edge labels}

Fix integers $R,S\ge0$ and set
\begin{equation}\label{eq:Q-smooth}
 Q=2^R3^S.
\end{equation}
For $t\in[0,1)$ define the vertex potentials
\begin{equation}\label{eq:vertex-potential}
 P_{a,b}(t)=\floor{2^a3^b t},
 \qquad 0\le a\le R,
 \quad 0\le b\le S.
\end{equation}
The horizontal and vertical edge labels are
\begin{align}
 h_{a,b}(t)
 &=P_{a+1,b}(t)-2P_{a,b}(t)
 =c_2(T_{2^a3^b}t)\in\{0,1\},
 \label{eq:h-carry}\\
 v_{a,b}(t)
 &=P_{a,b+1}(t)-3P_{a,b}(t)
 =c_3(T_{2^a3^b}t)\in\{0,1,2\}.
 \label{eq:v-carry}
\end{align}
Here $h_{a,b}$ is defined for $0\le a<R$, $0\le b\le S$, and $v_{a,b}$ for
$0\le a\le R$, $0\le b<S$.

\begin{proposition}[Plaquette compatibility]\label{prop:plaquette}
Every elementary square satisfies
\begin{equation}\label{eq:plaquette}
 3h_{a,b}+v_{a+1,b}=2v_{a,b}+h_{a,b+1}
 \qquad(0\le a<R,\ 0\le b<S).
\end{equation}
\end{proposition}

\begin{proof}
Compute the upper-right potential along the two directed paths:
\begin{align*}
 P_{a+1,b+1}
 &=3P_{a+1,b}+v_{a+1,b}
  =6P_{a,b}+3h_{a,b}+v_{a+1,b},\\
 P_{a+1,b+1}
 &=2P_{a,b+1}+h_{a,b+1}
  =6P_{a,b}+2v_{a,b}+h_{a,b+1}.
\end{align*}
Equating them gives \eqref{eq:plaquette}.
\end{proof}

\begin{definition}[Compatible carry rectangle]\label{def:compatible}
An $R\times S$ carry rectangle is a pair of arrays
\[
 h_{a,b}\in\{0,1\},
 \qquad
 v_{a,b}\in\{0,1,2\}
\]
on the edge sets above.  It is \emph{compatible} if every plaquette satisfies
\eqref{eq:plaquette}.
\end{definition}

\subsection{The exact cylinder--rectangle bijection}

Partition $[0,1)$ into the $Q$ half-open cylinders
\begin{equation}\label{eq:cylinders}
 I_k=\left[\frac{k}{Q},\frac{k+1}{Q}\right),
 \qquad 0\le k<Q.
\end{equation}
The next theorem is the finite combinatorial core of the report.

\begin{theorem}[Exact carry-rectangle classification]\label{thm:carry-bijection}
Let $R,S\ge0$ and $Q=2^R3^S$.

\begin{enumerate}
\item For every $0\le a\le R$, $0\le b\le S$, and every $t\in I_k$,
\begin{equation}\label{eq:potential-cylinder-formula}
 P_{a,b}(t)=\floor{\frac{k}{2^{R-a}3^{S-b}}}.
\end{equation}
Consequently the entire carry rectangle is constant on $I_k$.

\item Distinct cylinders produce distinct carry rectangles.

\item An abstract edge labeling is realized by some $t\in[0,1)$ if and only if it
is compatible in the sense of \Cref{def:compatible}.

\item Hence there are exactly
\begin{equation}\label{eq:number-compatible}
 2^R3^S=Q
\end{equation}
compatible $R\times S$ carry rectangles.
\end{enumerate}
\end{theorem}

\begin{proof}
Let $d=2^{R-a}3^{S-b}$.  If $t\in I_k$, then
\[
 \frac{k}{d}\le 2^a3^b t<\frac{k+1}{d}.
\]
No integer lies strictly between $k/d$ and $(k+1)/d$, and the right endpoint is
excluded.  Therefore the floor is $\floor{k/d}$, proving
\eqref{eq:potential-cylinder-formula}.

At the upper-right vertex, $P_{R,S}(t)=\floor{Qt}=k$.  Thus the rectangle recovers
$k$, so different cylinders give different rectangles.

Conversely, start with a compatible abstract labeling and set $P_{0,0}=0$.  Define
potentials along edges by
\begin{equation}\label{eq:potential-recursion}
 P_{a+1,b}=2P_{a,b}+h_{a,b},
 \qquad
 P_{a,b+1}=3P_{a,b}+v_{a,b}.
\end{equation}
The plaquette equation says exactly that the two recursions agree around every
square.  Since every two monotone lattice paths are related by adjacent square
interchanges, the potential is path-independent.  The digit ranges imply inductively
that
\begin{equation}\label{eq:potential-range}
 0\le P_{a,b}<2^a3^b.
\end{equation}
Put $k=P_{R,S}\in\{0,\ldots,Q-1\}$.  Running
\eqref{eq:potential-recursion} backward is Euclidean division by $2$ or $3$; it
recovers the unique quotients
\[
 P_{a,b}=\floor{\frac{k}{2^{R-a}3^{S-b}}}.
\]
By the first part, the labeling is therefore realized throughout $I_k$.
The bijection with the $Q$ cylinders proves the count.
\end{proof}

\begin{corollary}[Explicit edge digits]\label{cor:explicit-edge-digits}
For the rectangle associated with $k$, one has
\begin{align}
 h_{a,b}
 &=\floor{\frac{k}{2^{R-a-1}3^{S-b}}}
 -2\floor{\frac{k}{2^{R-a}3^{S-b}}},\label{eq:explicit-h}\\
 v_{a,b}
 &=\floor{\frac{k}{2^{R-a}3^{S-b-1}}}
 -3\floor{\frac{k}{2^{R-a}3^{S-b}}}.
 \label{eq:explicit-v}
\end{align}
Equivalently, these are the indicated quotient digits modulo $2$ and $3$.
\end{corollary}

\begin{proof}
Substitute \eqref{eq:potential-cylinder-formula} into
\eqref{eq:h-carry} and \eqref{eq:v-carry}.
\end{proof}

\subsection{Exact entropy collapse}

The edge alphabet contains $R(S+1)$ binary horizontal labels and $(R+1)S$ ternary
vertical labels.  Without compatibility there are
\begin{equation}\label{eq:naive-label-count}
 2^{R(S+1)}3^{(R+1)S}
\end{equation}
labelings.  Only $2^R3^S$ survive.

\begin{corollary}[Exact compatibility fraction and zero area entropy]
\label{cor:entropy-collapse}
The fraction of edge labelings that are compatible is
\begin{equation}\label{eq:compatibility-fraction}
 \frac{2^R3^S}{2^{R(S+1)}3^{(R+1)S}}=6^{-RS}.
\end{equation}
The logarithm of the number of feasible patterns is
\begin{equation}\label{eq:boundary-entropy}
 R\log2+S\log3,
\end{equation}
which is boundary-order rather than area-order.  In particular the two-dimensional
pattern entropy per plaquette tends to zero as $R,S\to\infty$.
\end{corollary}

This exact collapse is stronger than merely listing the plaquette equations.  It
shows that the restricted digit equations are independent in aggregate: each of the
$RS$ squares removes an exact factor $6$ from the naive count.

\begin{example}[The first nontrivial rectangle]\label{ex:one-square}
For $R=S=1$, $Q=6$.  Ordering the four edge labels as
$(h_{0,0},v_{0,0},h_{0,1},v_{1,0})$, the six compatible patterns are
\begin{center}
\begin{tabular}{c@{\qquad}c}
\toprule
$k$ & edge pattern \\
\midrule
$0$ & $(0,0,0,0)$\\
$1$ & $(0,0,1,1)$\\
$2$ & $(0,1,0,2)$\\
$3$ & $(1,1,1,0)$\\
$4$ & $(1,2,0,1)$\\
$5$ & $(1,2,1,2)$\\
\bottomrule
\end{tabular}
\end{center}
The plaquette equation is $3h_{0,0}+v_{1,0}=2v_{0,0}+h_{0,1}$.  Of the
$2^2 3^2=36$ possible edge labelings, exactly six satisfy it.
\end{example}

\subsection{Infinite compatible fields}

The finite theorem has a useful compact inverse-limit form.

\begin{theorem}[Global reconstruction from a compatible carry field]
\label{thm:global-reconstruction}
Suppose binary labels $h_{a,b}$ and ternary labels $v_{a,b}$ are given on the entire
first quadrant and satisfy every plaquette equation.  Reconstruct potentials from
$P_{0,0}=0$ by \eqref{eq:potential-recursion}.  Then the closed intervals
\begin{equation}\label{eq:nested-intervals}
 J_{a,b}=
 \left[\frac{P_{a,b}}{2^a3^b},
       \frac{P_{a,b}+1}{2^a3^b}\right]
\end{equation}
are compatible under refinement and have a unique common point $t\in[0,1]$.
For $t\in[0,1)$ the labels are the floor carries of $t$; the sole additional endpoint
field at $t=1$ is the all-maximal field $h\equiv1$, $v\equiv2$.
\end{theorem}

\begin{proof}
The recurrence says that a horizontal child interval is one of the two equal
subintervals of its parent, and a vertical child is one of the three equal
subintervals.  Plaquette compatibility makes these refinements independent of path.
Along the diagonal $(a,b)=(n,n)$, the interval length is $6^{-n}$, so the nested
closed intervals have a unique common point.  If $t<1$, the half-open version of the
finite cylinder formula identifies every floor and hence every edge label.  If the
intersection point is $1$, all potentials are $2^a3^b-1$, giving the maximal field.
\end{proof}

\begin{statusbox}{new finite and inverse-limit core}
\Cref{lem:carry-cocycle,prop:plaquette,thm:carry-bijection,cor:entropy-collapse,thm:global-reconstruction}
are \statusproved.  They require only Euclidean division, finite-grid path
independence, and nested intervals.  They are suitable for direct Lean
formalization without any transcendence theory.
\end{statusbox}

\section{Furstenberg universality and complete smooth residues}
\label{sec:furstenberg-universality}

\subsection{The external density input}

We use the following classical theorem in its special $2,3$ form.

\begin{theorem}[Furstenberg; Boshernitzan]\label{thm:furstenberg}
If $\alpha\in\R\setminus\Q$, then
\begin{equation}\label{eq:furstenberg-density}
 \bigl\{\fract{2^r3^s\alpha}:r,s\ge0\bigr\}
\end{equation}
is dense in $[0,1]$.
\end{theorem}

Furstenberg proved the result by topological dynamics
\cite{Furstenberg1967}; Boshernitzan later gave an elementary proof
\cite{Boshernitzan1994}.  The theorem is qualitative.  Quantitative variants are
discussed in \Cref{sec:hitting-times}.

A small but important tail observation will be used repeatedly.  For any fixed
$R_0,S_0$, the number $2^{R_0}3^{S_0}\alpha$ is still irrational, so
\begin{equation}\label{eq:tail-density}
 \bigl\{\fract{2^r3^s\alpha}:r\ge R_0,\ s\ge S_0\bigr\}
\end{equation}
is also dense.  Thus every open interval is visited arbitrarily far out in both
coordinates, not merely somewhere in the orbit.

\subsection{Every compatible finite carry pattern occurs}

\begin{theorem}[Universal finite carry language]\label{thm:universal-language}
Fix an irrational $\alpha$, dimensions $R,S$, and a compatible $R\times S$ carry
rectangle.  For every $R_0,S_0\ge0$, there exist $r\ge R_0$ and $s\ge S_0$ such that
the shifted carry block
\begin{equation}\label{eq:shifted-block}
 \bigl(h_{r+a,s+b}(\alpha),v_{r+a,s+b}(\alpha)\bigr)_{
 0\le a\le R,\ 0\le b\le S}
\end{equation}
is exactly the prescribed rectangle.
\end{theorem}

\begin{proof}
By \Cref{thm:carry-bijection}, the prescribed rectangle corresponds to a unique
cylinder $I_k$ of denominator $Q=2^R3^S$.  Its interior is nonempty.  Apply the tail
density \eqref{eq:tail-density} to choose $r\ge R_0,s\ge S_0$ with
$t=\fract{2^r3^s\alpha}$ in that interior.  The carries in the shifted block are the
carries of $t$ on the $R\times S$ rectangle, hence are the required pattern.
\end{proof}

\begin{corollary}[Finite-pattern no-go theorem]\label{cor:finite-pattern-nogo}
The set of finite compatible carry rectangles is independent of the irrational
number $\alpha$.  Consequently, no argument of the following form can prove the
Alaoglu--Erd\H{o}s conjecture:
\begin{quote}
associate to a hypothetical counterexample its $\times2,\times3$ carry field, inspect
only bounded windows of carry labels, and derive a contradiction by exhibiting a
compatible finite pattern that cannot occur.
\end{quote}
Every compatible bounded window already occurs for every irrational exponent.
\end{corollary}

This does not make the carry field useless.  It says precisely what additional data
must be retained: vertex heights, rational numerators, reduced denominator exponents,
prime support, or quantitative return times.  Discarding all of those and keeping
only a finite label window loses the arithmetic distinction.

\subsection{A complete residue theorem for smooth dilates}

The same cylinder bijection immediately yields an arithmetic consequence which seems
not to have been isolated in the supplied report.

\begin{theorem}[Complete floor residues modulo every $3$-smooth modulus]
\label{thm:complete-residues}
Let $\alpha\in\R\setminus\Q$ and $Q=2^R3^S$.  Then
\begin{equation}\label{eq:complete-residue-set}
 \left\{\floor{2^r3^s\alpha}\bmod Q:r\ge R,\ s\ge S\right\}
 =\Z/Q\Z.
\end{equation}
Moreover, every residue occurs with $r$ and $s$ arbitrarily large.
\end{theorem}

\begin{proof}
Fix $k\in\{0,\ldots,Q-1\}$.  By tail density, choose $r_0,s_0$ arbitrarily large
with
\[
 t=\fract{2^{r_0}3^{s_0}\alpha}\in I_k.
\]
Write $2^{r_0}3^{s_0}\alpha=n+t$.  Then
\begin{align*}
 \floor{2^{r_0+R}3^{s_0+S}\alpha}
 &=\floor{Q(n+t)}\\
 &=Qn+\floor{Qt}=Qn+k.
\end{align*}
Thus the floor is congruent to $k$ modulo $Q$.  Since $k$ was arbitrary, all residues
occur.  Choosing the anchor beyond any prescribed tail proves the final assertion.
\end{proof}

\begin{remark}[General two-base form]
The finite carry classification works verbatim for multiplicatively independent
integers $p,q\ge2$: horizontal digits lie in $\{0,\ldots,p-1\}$, vertical digits in
$\{0,\ldots,q-1\}$, the plaquette equation is
\[
 qh_{a,b}+v_{a+1,b}=pv_{a,b}+h_{a,b+1},
\]
and there are exactly $p^Rq^S$ compatible rectangles.  Furstenberg density then gives
complete residues modulo $p^Rq^S$.  The bases $2,3$ are special here only because of
the conjecture.
\end{remark}

\subsection{Why the theorem is stronger than ordinary equidistribution language}

The conclusion is not simply that the orbit is dense.  A single hit in $I_k$
encodes an entire coherent rectangle of $R(S+1)+(R+1)S$ edge labels and
$(R+1)(S+1)$ floor values.  Thus one real interval simultaneously prescribes all
lower-scale binary and ternary carries in the block.  The exact count
$2^R3^S$ explains why: the full rectangle contains no more information than its
upper-right floor residue.

This compression will be crucial below.  Under a hypothetical counterexample, the
same residue $k$ controls two rational systems at once.

\section{Arithmetic lift under a hypothetical counterexample}
\label{sec:counterexample-lift}

Throughout this section assume the inherited counterexample data
\eqref{eq:counterexample-output}--\eqref{eq:primitive-depth}.

\subsection{The exact block formula}

Fix an anchor $(r,s)$ and abbreviate
\begin{equation}\label{eq:anchor-abbrev}
 q=2^r3^s,
 \qquad
 \nu=\floor{q\beta},
 \qquad
 t=\fract{q\beta}.
\end{equation}
For local coordinates $a,b\ge0$, put $m_{a,b}=2^a3^b$.  Then
\begin{equation}\label{eq:block-floor-formula}
 \nu_{r+a,s+b}=m_{a,b}\nu+P_{a,b}(t),
\end{equation}
where $P_{a,b}(t)=\floor{m_{a,b}t}$ is the carry potential from
\eqref{eq:vertex-potential}.

\begin{proof}
Since $q\beta=\nu+t$,
\[
 \floor{m_{a,b}q\beta}
 =\floor{m_{a,b}\nu+m_{a,b}t}
 =m_{a,b}\nu+\floor{m_{a,b}t}.
\]
\end{proof}

The compact rational pair \eqref{eq:compact-pair} therefore obeys two synchronized
piecewise-monomial recurrences:
\begin{align}
 X_{r+a+1,s+b}
 &=\frac{X_{r+a,s+b}^{2}}{2^{h_{a,b}(t)}},
 &
 X_{r+a,s+b+1}
 &=\frac{X_{r+a,s+b}^{3}}{2^{v_{a,b}(t)}},
 \label{eq:X-recurrence}\\
 Y_{r+a+1,s+b}
 &=\frac{Y_{r+a,s+b}^{2}}{3^{h_{a,b}(t)}},
 &
 Y_{r+a,s+b+1}
 &=\frac{Y_{r+a,s+b}^{3}}{3^{v_{a,b}(t)}}.
 \label{eq:Y-recurrence}
\end{align}
The same binary and ternary edge digits drive both rational systems.  This is the
multiplicative two-dimensional analogue of the source report's one-dimensional
synchronized Sturmian recurrence.

\subsection{Reduced denominator exponents}

Recall the structural factorizations
\[
 M=2^\lambda W,\quad W\text{ odd},
 \qquad
 N=3^\mu Z,\quad 3\nmid Z.
\]
Define
\begin{equation}\label{eq:denominator-depths}
 d^{(2)}_{r,s}=\nu_{r,s}-\lambda q_{r,s},
 \qquad
 d^{(3)}_{r,s}=\nu_{r,s}-\mu q_{r,s}.
\end{equation}
Then \eqref{eq:compact-pair} is in lowest terms as
\begin{equation}\label{eq:reduced-pair}
 X_{r,s}=\frac{W^{q_{r,s}}}{2^{d^{(2)}_{r,s}}},
 \qquad
 Y_{r,s}=\frac{Z^{q_{r,s}}}{3^{d^{(3)}_{r,s}}}.
\end{equation}
The block formula gives
\begin{equation}\label{eq:denominator-block}
 d^{(j)}_{r+a,s+b}
 =m_{a,b}d^{(j)}_{r,s}+P_{a,b}(t),
 \qquad j\in\{2,3\}.
\end{equation}

\begin{theorem}[Synchronized complete denominator residues]
\label{thm:denominator-residues}
Let $Q=2^R3^S$.  For every residue $k\bmod Q$ and every prescribed tail
$r\ge r_0$, $s\ge s_0$, there is an anchor $(r,s)$ beyond that tail such that at the
top of the $R\times S$ block,
\begin{equation}\label{eq:denominator-congruences}
 d^{(2)}_{r+R,s+S}\equiv
 d^{(3)}_{r+R,s+S}\equiv k\pmod Q.
\end{equation}
Consequently each of the two reduced denominator exponent sequences attains every
residue modulo $Q$, and they do so at the same indices with the same residue.
\end{theorem}

\begin{proof}
Choose an anchor with $t_{r,s}\in I_k$ by
\Cref{thm:universal-language}.  At the top vertex,
$P_{R,S}(t)=\floor{Qt}=k$.  Since $m_{R,S}=Q$,
\eqref{eq:denominator-block} gives
\[
 d^{(j)}_{r+R,s+S}=Qd^{(j)}_{r,s}+k,
\]
which proves the congruences.
\end{proof}

\begin{corollary}[One full-depth branch]
Because $\min(\lambda,\mu)=0$, at least one of the two synchronized denominator
exponents is exactly the common floor $\nu_{r,s}$.  Thus one branch realizes the
complete residue theorem with no cancellation from a structural base factor.
\end{corollary}

\begin{statusbox}{arithmetic significance}
\Cref{thm:denominator-residues} is stronger than applying density separately to the
two compact rational coordinates.  It asserts synchronized residue universality for
two different reduced-denominator systems, with an inherited guarantee that at least
one system has maximal depth.  Nevertheless, the residue pattern itself is still
universal and therefore does not yet contradict integrality.
\end{statusbox}

\subsection{The two extreme rectangles}

The cylinders $I_0$ and $I_{Q-1}$ have especially rigid potentials.

\begin{proposition}[All-zero and all-maximal blocks]\label{prop:extreme-blocks}
Let $Q=2^R3^S$ and $m_{a,b}=2^a3^b$.

\begin{enumerate}
\item If $0<t<1/Q$, then
\begin{equation}\label{eq:zero-potential}
 P_{a,b}(t)=0,
 \qquad h_{a,b}=0,
 \qquad v_{a,b}=0
\end{equation}
throughout the rectangle, and
\begin{equation}\label{eq:zero-block-scaling}
 \nu_{r+a,s+b}=m_{a,b}\nu,
 \quad
 X_{r+a,s+b}=X_{r,s}^{m_{a,b}},
 \quad
 Y_{r+a,s+b}=Y_{r,s}^{m_{a,b}}.
\end{equation}

\item If $1-1/Q<t<1$, then
\begin{equation}\label{eq:max-potential}
 P_{a,b}(t)=m_{a,b}-1,
 \qquad h_{a,b}=1,
 \qquad v_{a,b}=2
\end{equation}
throughout the rectangle, and
\begin{equation}\label{eq:max-block-scaling}
 \nu_{r+a,s+b}=m_{a,b}\nu+m_{a,b}-1,
\end{equation}
\begin{equation}\label{eq:max-block-rationals}
 X_{r+a,s+b}=2\left(\frac{X_{r,s}}{2}\right)^{m_{a,b}},
 \qquad
 Y_{r+a,s+b}=3\left(\frac{Y_{r,s}}{3}\right)^{m_{a,b}}.
\end{equation}
\end{enumerate}
Both types of blocks occur arbitrarily far out.
\end{proposition}

\begin{proof}
For $0<t<1/Q$ and $m_{a,b}\le Q$, one has $0<m_{a,b}t<1$, so all potentials vanish.
For $1-1/Q<t<1$,
\[
 m_{a,b}-1<m_{a,b}t<m_{a,b},
\]
so $P_{a,b}=m_{a,b}-1$.  The carry labels and rational formulas follow from their
definitions.  Tail occurrence follows from Furstenberg density.
\end{proof}

The all-zero rectangles will be the quantitative focus.  They give exact powers over
an entire two-dimensional block, not merely one unusually small fractional part.


\section{Synchronized leading-digit gaps and residual arithmetic}
\label{sec:leading-gaps}

The all-zero carry rectangles admit a concrete interpretation in the ordinary
base-$2$ and base-$3$ expansions of the integral outputs.  This section records the
interpretation exactly and audits the most tempting attempts to extract a
contradiction from it.

\subsection{The paired residuals}

Fix an anchor $(r,s)$ and abbreviate
\begin{equation}\label{eq:anchor-abbreviations}
 q=2^r3^s,\qquad \nu=\floor{q\beta},\qquad t=\fract{q\beta}.
\end{equation}
Define the positive integral residuals
\begin{equation}\label{eq:paired-residuals}
 \Delta_2(q)=M^q-2^\nu,
 \qquad
 \Delta_3(q)=N^q-3^\nu.
\end{equation}
They satisfy the exact normalized identities
\begin{equation}\label{eq:normalized-residuals}
 \frac{\Delta_2(q)}{2^\nu}=2^t-1,
 \qquad
 \frac{\Delta_3(q)}{3^\nu}=3^t-1.
\end{equation}
In particular, when $0<t<1/Q$,
\begin{equation}\label{eq:residual-upper-bounds}
 0<\Delta_2(q)<\bigl(2^{1/Q}-1\bigr)2^\nu,
 \qquad
 0<\Delta_3(q)<\bigl(3^{1/Q}-1\bigr)3^\nu.
\end{equation}
The exponent $\nu$ is the same in both inequalities.  That synchronization is the
counterexample-specific input; a theorem about one base in isolation would not use
all the information.

\subsection{A simultaneous leading-zero theorem}

For $b\in\{2,3\}$ and $Q\ge2$, put
\begin{equation}\label{eq:gap-functions}
 G_b(Q)=\max\left\{0,
  \floor{-\log_b\bigl(b^{1/Q}-1\bigr)}\right\}.
\end{equation}
Thus
\begin{equation}\label{eq:gap-asymptotics}
 G_b(Q)=\log_b Q-\log_b(\log b)+O_b(Q^{-1})+O(1),
 \qquad Q\longrightarrow\infty,
\end{equation}
where the final $O(1)$ accounts for the integer part.

\begin{theorem}[Synchronized leading-zero gaps]\label{thm:leading-zero-gaps}
Assume the hypothetical counterexample normalization
\eqref{eq:counterexample-output}.  Let $Q=2^R3^S\ge2$.  There are arbitrarily large
$r$ and $s$ for which, with $q,\nu$ as in
\eqref{eq:anchor-abbreviations}, the following hold simultaneously.
\begin{enumerate}
\item The binary expansion of $M^q$ begins with a $1$ in position $\nu$, followed
by at least $G_2(Q)$ zero digits.
\item The ternary expansion of $N^q$ begins with a $1$ in position $\nu$, followed
by at least $G_3(Q)$ zero digits.
\item More explicitly,
\begin{equation}\label{eq:gap-expansions}
 M^q=2^\nu+\Delta_2(q),\quad 0<\Delta_2(q)<2^{\nu-G_2(Q)},
\end{equation}
\begin{equation}\label{eq:gap-expansions-three}
 N^q=3^\nu+\Delta_3(q),\quad 0<\Delta_3(q)<3^{\nu-G_3(Q)}.
\end{equation}
\end{enumerate}
\end{theorem}

\begin{proof}
By \Cref{prop:extreme-blocks}, choose arbitrarily remote anchors with
$0<t<1/Q$.  The first inequality in \eqref{eq:residual-upper-bounds} and the
definition of $G_2(Q)$ give
\[
 \Delta_2(q)<\bigl(2^{1/Q}-1\bigr)2^\nu\le 2^{\nu-G_2(Q)}.
\]
An integer smaller than $2^{\nu-G}$ cannot occupy any binary position from
$\nu-1$ down through $\nu-G$.  This proves the binary statement.  The ternary
argument is identical.  For $Q\ge2$ we also have $3^t<\sqrt3<2$, so the leading
ternary digit in position $\nu$ is indeed $1$, rather than $2$.
\end{proof}

\begin{corollary}[Arbitrarily remote paired gaps]
For every fixed $G\ge1$, a hypothetical counterexample would produce infinitely
many smooth indices $q=2^r3^s$ for which $M^q$ has at least $G$ binary zero digits
immediately after its leading digit and $N^q$ has at least $G$ ternary zero digits
immediately after its leading digit, with the leading positions equal.
\end{corollary}

\begin{proof}
Take $Q=2^R3^S$ large enough that both $G_2(Q)$ and $G_3(Q)$ are at least $G$, then
apply \Cref{thm:leading-zero-gaps}.
\end{proof}

\begin{statusbox}{what is genuinely new here}
Classical digit results can force the number of nonzero digits of powers to grow.
That does not forbid a long \emph{leading} gap along a sparse subsequence.  The new
constraint is paired: the binary and ternary gaps occur at the same smooth power and
have the same leading exponent $\nu$.  No theorem located in the literature search
converts precisely this synchronization into a contradiction.
\end{statusbox}

\subsection{Exact tail relation and the denominator tax}

Set
\begin{equation}\label{eq:u2-u3}
 u_2=\frac{\Delta_2(q)}{2^\nu},
 \qquad
 u_3=\frac{\Delta_3(q)}{3^\nu}.
\end{equation}
Then $1+u_2=2^t$ and $1+u_3=3^t$, so
\begin{equation}\label{eq:tail-power-relation}
 1+u_3=(1+u_2)^{\theta},
 \qquad
 \theta=\frac{\log3}{\log2},
\end{equation}
and equivalently
\begin{equation}\label{eq:tail-log-ratio}
 \frac{\log(1+u_3)}{\log(1+u_2)}=\theta.
\end{equation}
Both $u_2$ and $u_3$ are rational.  For an all-zero block,
$u_2,u_3=O(Q^{-1})$.  Taylor's theorem therefore gives
\begin{equation}\label{eq:rational-theta-approx}
 \left|
  \frac{u_3}{u_2}-\theta
 \right|\ll Q^{-1},
\end{equation}
with an absolute implied constant once $Q\ge2$.

At first sight \eqref{eq:rational-theta-approx} looks like a strong rational
approximation to $\theta$.  Written in integers, however,
\begin{equation}\label{eq:theta-approximant-integers}
 \frac{u_3}{u_2}
 =\frac{2^\nu\Delta_3(q)}{3^\nu\Delta_2(q)}.
\end{equation}
The numerator and denominator have height exponential in $q$, whereas the error is
only $O(Q^{-1})$.  Standard irrationality measures for $\log3/\log2$ therefore give
no contradiction unless $q$ is extraordinarily small as a function of $Q$.  This
is the same denominator tax that reappears in the hitting-time formulation below.

\subsection{Divisibility along an all-zero rectangle}

The exact scaling in \eqref{eq:zero-block-scaling} implies a family of honest
integral factorizations.  If $0<t<1/Q$ and $1\le m\le Q$, then
$\floor{mq\beta}=m\nu$ and
\begin{equation}\label{eq:residual-power-factorization}
 \Delta_2(mq)=M^{mq}-2^{m\nu}
 =\Delta_2(q)\sum_{j=0}^{m-1}M^{q(m-1-j)}2^{\nu j},
\end{equation}
\begin{equation}\label{eq:residual-power-factorization-three}
 \Delta_3(mq)=N^{mq}-3^{m\nu}
 =\Delta_3(q)\sum_{j=0}^{m-1}N^{q(m-1-j)}3^{\nu j}.
\end{equation}
Consequently
\begin{equation}\label{eq:residual-divisibility}
 \Delta_2(q)\mid\Delta_2(mq),
 \qquad
 \Delta_3(q)\mid\Delta_3(mq)
 \qquad(1\le m\le Q).
\end{equation}
One may refine these identities with the homogeneous cyclotomic factors
$\Phi_d(A,B)$:
\begin{equation}\label{eq:cyclotomic-residual}
 \Delta_2(mq)=\prod_{d\mid m}\Phi_d(M^q,2^\nu),
 \qquad
 \Delta_3(mq)=\prod_{d\mid m}\Phi_d(N^q,3^\nu).
\end{equation}

These are useful exact data, but they do not themselves improve the smallness
exponent.  Each quotient in \eqref{eq:residual-power-factorization} is naturally of
size about $m2^{(m-1)\nu}$ or $m3^{(m-1)\nu}$, exactly compensating for the growth of
the higher residual.  Primitive-divisor theorems add new primes to each branch but
do not couple the two branches.  A successful use of
\eqref{eq:residual-divisibility} therefore needs a genuinely joint invariant.

\section{Smooth hitting times and the quantitative bottleneck}
\label{sec:hitting-times}

The preceding conclusions are qualitative because Furstenberg's theorem does not
say how large the first multiplier must be.  The correct quantitative parameter is
as follows.

\begin{definition}[One-sided smooth hitting time]\label{def:hitting-time}
For irrational $\alpha$ and $Q\ge2$, define
\begin{equation}\label{eq:hitting-time}
 H_\alpha(Q)=\min\left\{2^r3^s:
 r,s\ge0,\quad 0<\fract{2^r3^s\alpha}<\frac1Q\right\}.
\end{equation}
Furstenberg density implies that $H_\alpha(Q)$ is finite.
\end{definition}

The one-sided definition targets the all-zero rectangle.  Replacing it by distance
to the nearest integer changes none of the scale comparisons below.

\subsection{A polynomial lower bound from logarithmic forms}

\begin{theorem}[Baker lower bound for a hypothetical counterexample]
\label{thm:baker-hitting-lower}
Assume \eqref{eq:counterexample-output}.  There exist effectively computable
constants $c_\beta>0$, $\eta_\beta>0$, and $Q_0(\beta)$ such that
\begin{equation}\label{eq:baker-hitting-lower}
 H_\beta(Q)\ge c_\beta Q^{\eta_\beta}
 \qquad(Q\ge Q_0(\beta)).
\end{equation}
\end{theorem}

\begin{proof}
Let $q=2^r3^s$ occur in the defining set of $H_\beta(Q)$ and put
$\nu=\floor{q\beta}$.  Since $M$ is not a power of $2$, the logarithms
$\log M$ and $\log2$ are linearly independent over $\Q$.  The nonzero linear form
\begin{equation}\label{eq:linear-form-hitting}
 \Lambda=q\log M-\nu\log2=\fract{q\beta}\log2
\end{equation}
therefore satisfies a standard effective lower bound
\begin{equation}\label{eq:linear-form-polynomial}
 |\Lambda|\ge c_1\max\{q,\nu,3\}^{-\kappa}
 \ge c_2q^{-\kappa}
\end{equation}
for constants $c_1,c_2>0$ and $\kappa>0$ depending only on $M$; see, for example,
Baker--W\"ustholz \cite{BakerWustholz1993}.  On the other hand,
$0<\Lambda<(\log2)/Q$.  Rearranging gives
$q\ge(c_2Q/\log2)^{1/\kappa}$, which has the stated form.
\end{proof}

The exponent in \eqref{eq:baker-hitting-lower} is not expected to be optimal.  Its
importance is qualitative: it excludes recurrence faster than every fixed power.

\subsection{The best available effective upper scale}

The ratio $\beta=\log M/\log2$ is Diophantine-generic by the same logarithmic-form
estimate.  It therefore lies in the scope of effective versions of Furstenberg's
theorem.  Bourgain--Lindenstrauss--Michel--Venkatesh proved an effective density
statement for Diophantine-generic points \cite{BLMV2009}.  Gayfulin and Moshchevitin
later gave an elementary explicit formulation whose density radius, in height
language, is essentially
\begin{equation}\label{eq:gm-density-radius}
 \bigl(\log\log\log X\bigr)^{-1/8+\varepsilon}
\end{equation}
for smooth multipliers at most a fixed power of $X$
\cite{GayfulinMoshchevitin2023}.

\begin{proposition}[Published effective upper scale]\label{prop:effective-upper}
Assume \eqref{eq:counterexample-output}.  For every $\varepsilon>0$ there is an
effective constant $C_{\beta,\varepsilon}>0$ such that, for all sufficiently large
$Q$,
\begin{equation}\label{eq:triple-exp-upper}
 H_\beta(Q)
 \le
 \exp\!\left(\exp\!\left(\exp\!\left(
 C_{\beta,\varepsilon}Q^{8+\varepsilon}
 \right)\right)\right).
\end{equation}
\end{proposition}

\begin{proof}
Apply the uniform effective density theorem of Gayfulin--Moshchevitin to a point in
the middle of $(0,1/Q)$, say $1/(2Q)$.  A density radius below $1/(2Q)$ guarantees
a smooth orbit point in that interval.  Solving
\[
 (\log\log\log X)^{-1/8+\varepsilon'}<\frac1{2Q}
\]
for $X$, and slightly enlarging the exponent to absorb $\varepsilon'$, yields the
triple-exponential expression.  Their theorem permits the smooth multiplier to be
bounded by a fixed power of $X$, which is absorbed by changing
$C_{\beta,\varepsilon}$.
\end{proof}

Combining the two estimates gives the present quantitative corridor
\begin{equation}\label{eq:hitting-corridor}
 c_\beta Q^{\eta_\beta}
 \ \le\ H_\beta(Q)\ \le\
 \exp\exp\exp\bigl(C_{\beta,\varepsilon}Q^{8+\varepsilon}\bigr).
\end{equation}
The gap between these scales is the principal obstruction exposed by the new route.
The exact constants are secondary; any subpolynomial upper bound would already be
far more than enough.

\subsection{A clean sufficient closing lemma}

\begin{problem}[Subpolynomial smooth recurrence target]
\label{target:subpolynomial-hitting}
Prove that if $M=2^\beta$ and $N=3^\beta$ are both integers, then along an unbounded
sequence of $3$-smooth moduli $Q$,
\begin{equation}\label{eq:subpolynomial-target}
 H_\beta(Q)=Q^{o(1)}.
\end{equation}
\end{problem}

\begin{theorem}[The target would prove Alaoglu--Erd\H{o}s]
\label{thm:target-closes}
If \Cref{target:subpolynomial-hitting} holds, then every real $x$ with
$2^x,3^x\in\Z$ is an integer.
\end{theorem}

\begin{proof}
Assume a nonintegral solution and pass to the primitive $\beta$.  The target gives,
for every $\delta>0$ and infinitely many $Q$, a smooth $q<Q^\delta$ with
$0<\{q\beta\}<1/Q$.  But \Cref{thm:baker-hitting-lower} gives
$q\ge c_\beta Q^{\eta_\beta}$.  Taking $\delta<\eta_\beta$ and $Q$ sufficiently
large is impossible.
\end{proof}

\begin{statusbox}{the exact new wall}
Topological dynamics supplies every finite rectangle but no useful first-occurrence
bound.  Transcendence theory supplies a polynomial exclusion zone but no recurrence.
The missing theorem must use simultaneous integrality to improve the recurrence
scale.  Improving the constants in either existing estimate without crossing the
polynomial/subpolynomial boundary will not settle the conjecture by this route.
\end{statusbox}

\section{Metric rarity and the topological--arithmetic boundary}
\label{sec:metric-rarity}

The target \eqref{eq:subpolynomial-target} is not a generic dynamical property.  It
would place a hypothetical counterexample in an exceptionally recurrent class which
is invisible both to finite carry languages and to ordinary metric theory.

\begin{definition}[Smooth approximation exponent]
For irrational $\alpha$, define
\begin{equation}\label{eq:smooth-omega}
 \omega_{2,3}(\alpha)=
 \limsup_{r+s\to\infty}
 \frac{-\log\norm{2^r3^s\alpha}_{\R/\Z}}
      {r\log2+s\log3},
\end{equation}
where $\norm{y}_{\R/\Z}$ is distance to the nearest integer.
\end{definition}

\begin{proposition}[Baker finiteness]
If $\alpha=\log A/\log B$ for multiplicatively independent positive integers
$A,B>1$, then
\begin{equation}\label{eq:omega-finite}
 \omega_{2,3}(\alpha)<\infty.
\end{equation}
In particular this holds for the hypothetical $\beta$.
\end{proposition}

\begin{proof}
For $q=2^r3^s$, a nearest integer $p$ gives the nonzero logarithmic form
$q\log A-p\log B$.  A polynomial lower bound in $q$ implies
$\norm{q\alpha}_{\R/\Z}\gg q^{-\kappa}$ for a fixed $\kappa$, and hence the ratio
in \eqref{eq:smooth-omega} is bounded above by $\kappa+o(1)$.
\end{proof}

\begin{theorem}[Metric value]\label{thm:metric-omega}
For Lebesgue-almost every $\alpha\in[0,1]$,
\begin{equation}\label{eq:metric-omega-zero}
 \omega_{2,3}(\alpha)=0.
\end{equation}
\end{theorem}

\begin{proof}
Fix $\sigma>0$ and let
\[
 E_{r,s}(\sigma)=
 \left\{\alpha\in[0,1]:
 \norm{2^r3^s\alpha}_{\R/\Z}<(2^r3^s)^{-\sigma}
 \right\}.
\]
Multiplication by an integer preserves Lebesgue measure on the circle, so
\[
 |E_{r,s}(\sigma)|\le2(2^r3^s)^{-\sigma}.
\]
The double sum factors and converges:
\[
 \sum_{r,s\ge0}|E_{r,s}(\sigma)|
 \le2\left(\sum_{r\ge0}2^{-\sigma r}\right)
       \left(\sum_{s\ge0}3^{-\sigma s}\right)<\infty.
\]
The first Borel--Cantelli lemma therefore says that almost every $\alpha$ belongs to
only finitely many $E_{r,s}(\sigma)$.  Intersecting over positive rational $\sigma$
gives $\omega_{2,3}(\alpha)\le0$.  The definition gives the reverse inequality.
\end{proof}

\begin{proposition}[Subpolynomial hitting is infinite-exponent recurrence]
If \eqref{eq:subpolynomial-target} holds along an unbounded sequence, then
\begin{equation}\label{eq:omega-infinite}
 \omega_{2,3}(\beta)=+\infty.
\end{equation}
\end{proposition}

\begin{proof}
Choose $Q_j\to\infty$ and smooth $q_j\le Q_j^{\varepsilon_j}$, with
$\varepsilon_j\to0$, such that $0<\{q_j\beta\}<1/Q_j$.  Then
\[
 \frac{-\log\norm{q_j\beta}_{\R/\Z}}{\log q_j}
 \ge\frac{\log Q_j}{\varepsilon_j\log Q_j}
 =\varepsilon_j^{-1}\longrightarrow\infty.
\]
\end{proof}

Thus the closing target asks simultaneous integrality to force a recurrence behavior
which is not merely measure zero but incompatible with every finite irrationality
measure.  This explains why a proof cannot come from Furstenberg minimality alone:
minimality is shared by every irrational point, whereas the desired recurrence is at
the opposite extreme from the metric norm.

\subsection{A hierarchy of information}

The new approach separates four levels which are easy to conflate.
\begin{enumerate}
\item \textbf{Finite compatibility:} the plaquette equation and digit ranges.
This is completely classified by \Cref{thm:carry-bijection}.
\item \textbf{Topological recurrence:} every compatible finite pattern occurs.
This holds for every irrational point by \Cref{thm:universal-language}.
\item \textbf{Quantitative recurrence:} how soon a specified pattern occurs.
This is measured by $H_\alpha(Q)$ and is not controlled by the finite language.
\item \textbf{Arithmetic acceleration:} an improvement caused by the simultaneous
rational representations $2^\beta=M$ and $3^\beta=N$.  This is the missing level.
\end{enumerate}
Any argument which uses only levels 1 and 2 is ruled out by the finite-pattern no-go
theorem.  Any argument which reaches level 3 but retains the published
triple-exponential scale remains far from the Baker boundary.

\section{Digit lifts, multiplication carries, and integral defect sequences}
\label{sec:digit-lifts}

A second route compares the ordinary digits of the two integer outputs.  It is
independent of the two-dimensional floor-carry field, but the two routes meet on the
all-zero rectangles.

\subsection{The two digit-lift maps}

Write a nonnegative integer $m$ in binary and a nonnegative integer $n$ in ternary:
\begin{equation}\label{eq:digit-expansions}
 m=\sum_{j\ge0}\varepsilon_j2^j,
 \quad \varepsilon_j\in\{0,1\},
 \qquad
 n=\sum_{j\ge0}d_j3^j,
 \quad d_j\in\{0,1,2\}.
\end{equation}
Define
\begin{equation}\label{eq:digit-lift-definitions}
 L(m)=\sum_{j\ge0}\varepsilon_j3^j,
 \qquad
 R(n)=\sum_{j\ge0}d_j2^j.
\end{equation}
Thus $L$ reads a binary digit string in base $3$, while $R$ reads a ternary digit
string at the value $2$.  Put
\begin{equation}\label{eq:theta-rho}
 \theta=\log_2 3>1,
 \qquad
 \rho=\log_3 2=\theta^{-1}\in(0,1).
\end{equation}

\begin{theorem}[Sharp direction of the digit lifts]\label{thm:digit-lift-power}
The following statements hold.
\begin{enumerate}
\item $L$ is strictly increasing on $\N\).
\item For every $m\ge0$,
\begin{equation}\label{eq:L-power-bound}
 L(m)\le m^\theta,
\end{equation}
with equality exactly when $m=0$ or $m$ is a power of $2$.
\item More quantitatively, if $m=2^B+r$ with $0<r<2^B$, then
\begin{equation}\label{eq:L-quantitative-defect}
 m^\theta-L(m)\ge
 \left(\frac32\right)^B r.
\end{equation}
\item For every $n\ge0$,
\begin{equation}\label{eq:R-power-bound}
 n^\rho\le R(n),
\end{equation}
with equality exactly when $n=0$ or $n$ is a power of $3$.
\end{enumerate}
\end{theorem}

\begin{proof}
For monotonicity, let $B$ be the highest binary position where $m<n$ differ.  The
contribution $3^B$ of the new leading digit exceeds the largest possible loss from
all lower positions:
\[
 3^B-\sum_{j=0}^{B-1}3^j=\frac{3^B+1}{2}>0.
\]

We prove the quantitative estimate by induction on $B$.  The scalar inequality
\begin{equation}\label{eq:scalar-convexity}
 (1+u)^\theta\ge1+u^\theta+u
 \qquad(0\le u\le1)
\end{equation}
holds because the difference has value $0$ at both endpoints and has negative
second derivative on $(0,1)$; here $2^\theta=3$ is used at $u=1$.  With
$a=2^B$ and $u=r/a$, induction gives $L(r)\le r^\theta$, and hence
\[
 \begin{split}
 m^\theta-L(m)
 &=(a+r)^\theta-a^\theta-L(r)\\
 &\ge(a+r)^\theta-a^\theta-r^\theta\\
 &\ge a^{\theta-1}r
 =\left(\frac32\right)^B r.
 \end{split}
\]
This is strict for $r>0$, while repeated equality occurs precisely for powers of
$2$.

For $R$, use subadditivity of $x\mapsto x^\rho$:
\[
 n^\rho
 =\left(\sum_jd_j3^j\right)^\rho
 \le\sum_j(d_j3^j)^\rho
 =\sum_jd_j^\rho2^j
 \le\sum_jd_j2^j=R(n).
\]
Equality requires one nonzero summand and that its digit be $1$, which gives a power
of $3$.
\end{proof}

\subsection{Composition and exact valuations}

\begin{proposition}[Composition laws]\label{prop:digit-composition}
For every $m,n\ge0$,
\begin{equation}\label{eq:R-L-identity}
 R(L(m))=m,
\end{equation}
while
\begin{equation}\label{eq:L-R-inflation}
 L(R(n))\ge n.
\end{equation}
Equality in \eqref{eq:L-R-inflation} holds exactly when the ternary normalization of
$n$ requires no binary carries after its digits are read at $2$; in particular it
holds whenever all ternary digits are $0$ or $1$.
\end{proposition}

\begin{proof}
The ternary expansion of $L(m)$ consists of the binary digits of $m$, all belonging
to $\{0,1\}$, so reading it back at $2$ gives $m$.  For the other direction, start
with the ternary digit polynomial $\sum d_jX^j$.  Its value at $2$ is $R(n)$.
Normalizing the coefficients in base $2$ replaces two units in position $j$ by one
unit in position $j+1$.  When evaluated at $3$, each such carry changes
$2\cdot3^j$ into $3^{j+1}$ and therefore increases the value by $3^j$.  The initial
value at $3$ is $n$ and the normalized value is $L(R(n))$.
\end{proof}

\begin{proposition}[Valuation transfer]\label{prop:digit-valuation}
For every positive integer $m$,
\begin{equation}\label{eq:L-valuation}
 v_3(L(m))=v_2(m).
\end{equation}
For every positive integer $n$,
\begin{equation}\label{eq:R-valuation}
 v_2(R(n))\ge v_3(n),
\end{equation}
with equality if the least nonzero ternary digit of $n$ is $1$.
\end{proposition}

\begin{proof}
If the least binary $1$ of $m$ occurs in position $k$, then
$L(m)=3^k(1+3u)$ for an integer $u$.  If the least nonzero ternary digit of $n$ is
$d_k$, then $R(n)=2^k(d_k+2u)$; the parenthesis is odd when $d_k=1$ and is even when
$d_k=2$.
\end{proof}

\subsection{Multiplication carries}

The maps are not ring homomorphisms, but their failure has a definite sign and an
exact carry expansion.

\begin{theorem}[Signed multiplicativity and carry gains]
\label{thm:signed-multiplicativity}
For all nonnegative integers $a,b$,
\begin{equation}\label{eq:L-supermultiplicative}
 L(ab)\ge L(a)L(b),
\end{equation}
\begin{equation}\label{eq:R-submultiplicative}
 R(ab)\le R(a)R(b).
\end{equation}
More precisely, there are finitely supported nonnegative integer sequences
$\kappa_j(a,b)$ and $\lambda_j(a,b)$ such that
\begin{equation}\label{eq:binary-carry-gain}
 L(ab)-L(a)L(b)=\sum_{j\ge0}\kappa_j(a,b)3^j,
\end{equation}
\begin{equation}\label{eq:ternary-carry-loss}
 R(a)R(b)-R(ab)=\sum_{j\ge0}\lambda_j(a,b)2^j.
\end{equation}
The coefficients count elementary carries in a normalization algorithm.
\end{theorem}

\begin{proof}
Multiply the binary digit polynomials of $a$ and $b$.  Before normalization, their
value at $3$ is $L(a)L(b)$.  Every elementary base-$2$ carry replaces two units at
position $j$ by one at position $j+1$ and increases the value at $3$ by exactly
$3^j$.  After all carries the polynomial is the binary expansion of $ab$, whose
value at $3$ is $L(ab)$.  Summing the carry gains proves
\eqref{eq:binary-carry-gain}.  The ternary case is analogous: replacing three units
at position $j$ by one at $j+1$ decreases evaluation at $2$ by exactly $2^j$.
\end{proof}

\begin{remark}
The strictness conditions are concrete.  Equality in \eqref{eq:L-supermultiplicative}
means that the convolution of the two binary digit strings needs no carry; equality
in \eqref{eq:R-submultiplicative} means the corresponding ternary convolution needs
no carry.  This packages multiplication of output powers into another nonnegative
integer cocycle.
\end{remark}

\subsection{Consequences for a hypothetical output pair}

For $n\ge1$, define
\begin{equation}\label{eq:power-defects}
 D_n=N^n-L(M^n),
 \qquad
 E_n=R(N^n)-M^n.
\end{equation}

\begin{theorem}[Positive integral digit defects]\label{thm:positive-defects}
Under \eqref{eq:counterexample-output}, $D_n,E_n$ are positive integers for every
$n\ge1$.  Moreover,
\begin{equation}\label{eq:defect-growth-brackets}
 \frac{N^n}{3}\le L(M^n)<N^n,
 \qquad
 M^n<R(N^n)<4M^n,
\end{equation}
and hence
\begin{equation}\label{eq:digit-growth-rates}
 \lim_{n\to\infty}L(M^n)^{1/n}=N,
 \qquad
 \lim_{n\to\infty}R(N^n)^{1/n}=M.
\end{equation}
\end{theorem}

\begin{proof}
Because $M$ is not a power of $2$, neither is $M^n$, so
\Cref{thm:digit-lift-power} gives
$L(M^n)<M^{n\theta}=N^n$.  The reverse inequality is strict because $N^n$ is not a
power of $3$.  If $B=\floor{n\beta}$, then the leading binary digit of $M^n$ is in
position $B$, so $L(M^n)\ge3^B=N^n/3^{\{n\beta\}}\ge N^n/3$.  The ternary expansion
of $N^n<3^{B+1}$ uses positions at most $B$, whence
$R(N^n)\le2\sum_{j=0}^{B}2^j<2^{B+2}\le4M^n$.  Taking $n$th roots proves the limits.
\end{proof}

Thus the digit-lift route cannot win by a fixed exponential gap: the two lifted
sequences have exactly the conjectural exponential rates.  The useful information is
in their subexponential defects.

Let
\begin{equation}\label{eq:power-tail-notation}
 B_n=\floor{n\beta},\quad t_n=\fract{n\beta},\quad
 r_n=M^n-2^{B_n},\quad s_n=N^n-3^{B_n}.
\end{equation}
Then $0<r_n<2^{B_n}$ and $0<s_n<2\cdot3^{B_n}$.

\begin{theorem}[Tail inequalities and a sharp defect sandwich]
\label{thm:tail-defect-sandwich}
For every $n\ge1$,
\begin{equation}\label{eq:L-tail-inequality}
 L(r_n)<s_n,
\end{equation}
and
\begin{equation}\label{eq:D-sandwich}
 3^{B_n}\bigl(2^{t_n}-1\bigr)
 \le D_n
 \le3^{B_n}\bigl(3^{t_n}-1\bigr).
\end{equation}
If $t_n<\rho=\log_3 2$, then the leading ternary digit of $N^n$ is $1$ and
\begin{equation}\label{eq:R-tail-inequality}
 r_n<R(s_n).
\end{equation}
In that range,
\begin{equation}\label{eq:E-sandwich}
 (1-\rho)2^{B_n}\bigl(3^{t_n}-1\bigr)^\rho
 \le E_n
 <4\,2^{B_n}\bigl(3^{t_n}-1\bigr)^\rho.
\end{equation}
\end{theorem}

\begin{proof}
Since the leading binary digit of $M^n$ is isolated,
$L(M^n)=3^{B_n}+L(r_n)$.  Subtracting from
$N^n=3^{B_n}+s_n$ gives $D_n=s_n-L(r_n)>0$, proving
\eqref{eq:L-tail-inequality}.  The quantitative lower bound
\eqref{eq:L-quantitative-defect}, with
$r_n=2^{B_n}(2^{t_n}-1)$, gives the lower half of
\eqref{eq:D-sandwich}.  The upper half follows from
$L(M^n)\ge3^{B_n}$.

If $t_n<\rho$, then $N^n<2\cdot3^{B_n}$, so its leading ternary digit is $1$ and
$R(N^n)=2^{B_n}+R(s_n)$.  Hence
$E_n=R(s_n)-r_n>0$.  Put $u=3^{t_n}-1\in(0,1)$.  By
\eqref{eq:R-power-bound}, concavity of $x^\rho$, and $u\le u^\rho$,
\[
 \begin{split}
 E_n
 &\ge s_n^\rho-2^{B_n}\bigl((1+u)^\rho-1\bigr)\\
 &\ge2^{B_n}(u^\rho-\rho u)
 \ge(1-\rho)2^{B_n}u^\rho.
 \end{split}
\]
For the upper bound, $E_n<R(s_n)$.  If the highest ternary position of $s_n$ is
$J$, then
$R(s_n)\le2\sum_{j=0}^{J}2^j<4\cdot2^J\le4s_n^\rho$, which gives
\eqref{eq:E-sandwich}.
\end{proof}

For an all-zero smooth anchor $n=q$ with $t_q<1/Q$, the theorem gives the paired
relative estimates
\begin{equation}\label{eq:defect-smallness}
 \frac{D_q}{N^q}=\Theta(t_q)=O(Q^{-1}),
 \qquad
 \frac{E_q}{M^q}=\Theta\bigl(t_q^\rho\bigr)=O(Q^{-\rho}),
\end{equation}
with effective constants independent of the anchor.  The differing exponents reflect
the H\"older distortion between ternary and binary digit evaluation.

\begin{proposition}[Defect valuation transfer]\label{prop:defect-valuation}
Let $\lambda=v_2(M)$ and $\mu=v_3(N)$.  If $\lambda\ne\mu$, then for every $n\ge1$,
\begin{equation}\label{eq:D-valuation}
 v_3(D_n)=n\min\{\lambda,\mu\}.
\end{equation}
In particular, under the primitive condition \eqref{eq:primitive-depth}, if exactly
one of $\lambda,\mu$ is positive, then every $D_n$ is a $3$-adic unit.
\end{proposition}

\begin{proof}
By \Cref{prop:digit-valuation},
$v_3(L(M^n))=v_2(M^n)=n\lambda$, while
$v_3(N^n)=n\mu$.  When these valuations differ, the valuation of their difference is
the smaller one.
\end{proof}

\subsection{An exact bridge from digit carries to analytic defects}

Define the nonnegative carry functionals
\begin{equation}\label{eq:carry-functionals}
 C_2(a,b)=L(ab)-L(a)L(b),
 \qquad
 C_3(a,b)=R(a)R(b)-R(ab).
\end{equation}

\begin{theorem}[Defect--carry recurrences]\label{thm:defect-carry-recurrence}
For all positive integers $m,n$,
\begin{equation}\label{eq:D-carry-recurrence}
 C_2(M^m,M^n)
 =N^mD_n+N^nD_m-D_mD_n-D_{m+n},
\end{equation}
\begin{equation}\label{eq:E-carry-recurrence}
 C_3(N^m,N^n)
 =M^mE_n+M^nE_m+E_mE_n-E_{m+n}.
\end{equation}
\end{theorem}

\begin{proof}
Write $L(M^k)=N^k-D_k$ in
$C_2(M^m,M^n)=L(M^{m+n})-L(M^m)L(M^n)$ and expand.  The second identity follows by
writing $R(N^k)=M^k+E_k$.
\end{proof}

\begin{corollary}[Carry sparsity inside an all-zero block]
\label{cor:carry-sparsity}
Suppose $0<t_q<1/Q$ and $m+n\le Q$.  Then
\begin{equation}\label{eq:C2-small}
 0\le C_2(M^{mq},M^{nq})
 \ll \frac{m+n}{Q}\,N^{(m+n)q},
\end{equation}
\begin{equation}\label{eq:C3-small}
 0\le C_3(N^{mq},N^{nq})
 \ll \left(\frac{m}{Q}\right)^\rho M^{(m+n)q}
     +\left(\frac{n}{Q}\right)^\rho M^{(m+n)q},
\end{equation}
with absolute implied constants.
\end{corollary}

\begin{proof}
For $k\le Q$, one has $t_{kq}=kt_q<1$.  Apply
\Cref{thm:tail-defect-sandwich} to bound
$D_{kq}/N^{kq}\ll kt_q\ll k/Q$ and, when $kt_q<\rho$,
$E_{kq}/M^{kq}\ll(kt_q)^\rho$.  If one of $m,n,m+n$ lies outside that smaller
range, enlarge the absolute constant using \eqref{eq:defect-growth-brackets}.
Substitution in \Cref{thm:defect-carry-recurrence} gives the result.
\end{proof}

The recurrence is a concrete point of contact between the two new packages.  A
hypothetical counterexample forces multiplication of enormous powers to incur only a
small high-base carry gain along the same smooth subsequences that create the floor
carry rectangles.  At present, general lower bounds for these multiplication-carry
functionals are too weak: the first unavoidable carry may occur far below the leading
position, exactly as permitted by the leading-zero gaps.

\section{Stress tests: what the new information does not prove}
\label{sec:stress-tests}

A useful progress report should not merely accumulate true statements; it should
identify which combinations have actually been tested and why they stop.  The table
below records the principal audits.

\begin{longtable}{>{\raggedright\arraybackslash}p{0.24\textwidth}>{\raggedright\arraybackslash}p{0.31\textwidth}>{\raggedright\arraybackslash}p{0.35\textwidth}}
\toprule
Attempt & Exact input & Failure point or surviving target \\
\midrule
\endfirsthead
\toprule
Attempt & Exact input & Failure point or surviving target \\
\midrule
\endhead
Forbid a finite carry pattern & Plaquette compatibility and the counterexample
carry field & Every compatible finite pattern occurs for every irrational exponent.
This strategy is ruled out by \Cref{cor:finite-pattern-nogo}.\\

Use the rarity of an all-zero rectangle & Arbitrarily long simultaneous leading gaps
from \Cref{thm:leading-zero-gaps} & Such rectangles also occur for every irrational
orbit.  The missing information is the first hitting scale, not existence.\\

Approximate $\log3/\log2$ by the residual ratio &
Equation \eqref{eq:rational-theta-approx} & The rational denominator is exponential
in the smooth multiplier $q$.  An $O(Q^{-1})$ error is far too weak unless
$q=Q^{o(1)}$.\\

Exploit ordinary cyclotomic factors &
Equations \eqref{eq:residual-power-factorization}--\eqref{eq:cyclotomic-residual} &
The quotient factors have exactly the size needed to absorb the residual growth.
Primitive primes occur separately in the two branches and are not synchronized.\\

Use a theorem on sparse digits of powers & Paired leading gaps & Known single-base
sparsity results force many nonzero digits globally, not the absence of an arbitrarily
long leading gap on a sparse subsequence.  A joint theorem with the same leading
position would be needed.\\

Use supermultiplicativity of $L$ &
$L(M^{m+n})\ge L(M^m)L(M^n)$ & The exponential growth rate is exactly $N$ by
\eqref{eq:digit-growth-rates}; there is no fixed exponential deficit to amplify.\\

Use the carry gains $C_2,C_3$ &
The exact recurrences \eqref{eq:D-carry-recurrence}--\eqref{eq:E-carry-recurrence} &
The first mandatory multiplication carry can lie far below the leading digit.
Current lower bounds match the scale allowed by the leading gaps.\\

Use complete residues modulo smooth $Q$ &
\Cref{thm:denominator-residues} & Residue universality is itself a consequence of
density and is compatible with every irrational point.  One needs a height,
valuation, or factorization consequence attached to a selected residue.\\

Replace density by equidistribution & The orbit $\{2^r3^s\beta\}$ & Furstenberg's
theorem gives density, not pointwise equidistribution for every irrational.  Assuming
normality or uniform discrepancy would insert an unproved theorem.\\

Apply the effective $\times2,\times3$ theorem directly &
\Cref{prop:effective-upper} & The available upper scale is roughly triple
exponential in a power of $Q$, while Baker excludes only a polynomial scale.\\

Infer a theorem from the reported rival breakthrough & Current announcements & No
public Alaoglu--Erd\H{o}s statement, proof object, or Lean source was found.  The
claim cannot be used as a lemma.\\
\bottomrule
\end{longtable}

\subsection{Three false shortcuts in formula form}

\paragraph{Density is not a rate.}
The implication
\[
 \overline{\{2^r3^s\beta:r,s\ge0\}}=\T
 \quad\Longrightarrow\quad
 H_\beta(Q)\le Q^C
\]
is false as a matter of logic.  Compactness proves finiteness for each fixed $Q$ but
supplies no uniform dependence on $Q$.

\paragraph{A small real residual need not have a large structural valuation.}
When $M$ is odd, $\Delta_2(q)=M^q-2^\nu$ is odd, even when it is tiny relative to
$2^\nu$.  Thus archimedean smallness does not automatically create $2$-adic depth.
The analogous statement holds for the ternary branch.

\paragraph{Two small forms are not independent forms.}
The expressions
\[
 q\log M-\nu\log2=t\log2,
 \qquad
 q\log N-\nu\log3=t\log3
\]
are proportional because of the original logarithmic relation.  Multiplying their
individual Baker lower bounds does not create a two-dimensional determinant; doing
so would disguise the four-exponentials barrier rather than cross it.

\section{Computational reconnaissance}
\label{sec:computation}

The finite classification is exact and does not require computation.  Computation is
still useful for checking indexing conventions, exploring first-occurrence scales,
and preparing formal tests.  All calculations in this section are labeled
\statuscomputed{} and are logically dispensable.

\subsection{Exact enumeration of compatible rectangles}

For a fixed $R,S$, the following integer-only pseudocode constructs every compatible
rectangle.  It is a direct executable version of \Cref{thm:carry-bijection}.

\begin{lstlisting}[language=Python,caption={Exact enumeration from the cylinder index}]
def carry_rectangle(R: int, S: int, k: int):
    Q = (2**R) * (3**S)
    assert 0 <= k < Q

    P = {}
    for a in range(R + 1):
        for b in range(S + 1):
            denominator = (2 ** (R - a)) * (3 ** (S - b))
            P[a, b] = k // denominator

    h = {(a, b): P[a + 1, b] - 2 * P[a, b]
         for a in range(R) for b in range(S + 1)}
    v = {(a, b): P[a, b + 1] - 3 * P[a, b]
         for a in range(R + 1) for b in range(S)}

    assert all(x in (0, 1) for x in h.values())
    assert all(x in (0, 1, 2) for x in v.values())
    assert all(3*h[a, b] + v[a + 1, b]
               == 2*v[a, b] + h[a, b + 1]
               for a in range(R) for b in range(S))
    return P, h, v
\end{lstlisting}

Running over $0\le k<2^R3^S$ produces exactly that many distinct pairs $(h,v)$.
No floating-point operation occurs.

\subsection{First complete cylinder coverage in bounded boxes}

A separate high-precision scan asked when a bounded orbit box had hit every cylinder
$I_k$ for $Q=6,36,216$.  The exponents were restricted to $0\le r,s\le60$ and
$120$ decimal digits were used.  The quantity reported is the largest, over all
cylinders, of the logarithm of the first multiplier $2^r3^s$ found for that cylinder.

\begin{center}
\begin{tabular}{lrrr}
\toprule
$\alpha$ & $Q=6$ & $Q=36$ & $Q=216$\\
\midrule
$\log_2 3$ & $4.852030264$ & $21.48756260$ & $40.80500268$\\
$\log_2 5$ & $3.871201011$ & $15.77248612$ & $42.76666119$\\
$\sqrt2$   & $4.852030264$ & $19.70935416$ & $52.07880764$\\
$\pi$      & $2.484906650$ & $9.416378455$ & $40.56943661$\\
\bottomrule
\end{tabular}
\end{center}

All cylinders were found within the scan box in these twelve cases.  The data show
large fluctuations and no stable power law at these tiny scales.  In particular,
they do not support extrapolation to the target
\eqref{eq:subpolynomial-target}.  Their purpose is only to validate the cylinder
indexing and to illustrate that first complete coverage can be controlled by a very
different orbit point for nearby examples.

\subsection{Reproducible diagnostic script}

\begin{lstlisting}[language=Python,caption={High-precision bounded orbit scan}]
import mpmath as mp

mp.mp.dps = 120
alphas = {
    "log_2(3)": mp.log(3) / mp.log(2),
    "log_2(5)": mp.log(5) / mp.log(2),
    "sqrt(2)": mp.sqrt(2),
    "pi": mp.pi,
}

for R, S in [(1, 1), (2, 2), (3, 3)]:
    Q = (2**R) * (3**S)
    print("Q", Q)
    for name, alpha in alphas.items():
        first = [None] * Q
        for r in range(61):
            for s in range(61):
                q = (2**r) * (3**s)
                t = mp.frac(mp.mpf(q) * alpha)
                k = int(mp.floor(Q * t))
                candidate = (mp.log(q), r + s, r, s)
                if first[k] is None or candidate[0] < first[k][0]:
                    first[k] = candidate
        assert all(item is not None for item in first)
        worst = max(first, key=lambda item: item[0])
        print(name, mp.nstr(worst[0], 12),
              max(item[1] for item in first))
\end{lstlisting}

A proof-oriented implementation should replace transcendental floating-point
comparisons by interval enclosures and record the distance from every sampled point
to the nearest cylinder boundary.  None of the theorem proofs above depend on this
scan.

\section{Lean formalization blueprint}
\label{sec:lean-blueprint}

The finite carry theory is unusually suitable for formalization because its core can
be expressed entirely with natural-number division and remainders.  The analytic and
external dynamical inputs can be attached later.

\subsection{Proposed module graph}

\begin{center}
\begin{tabularx}{0.96\textwidth}{>{\ttfamily\raggedright\arraybackslash}p{0.27\textwidth}X}
\toprule
Module & Principal contents \\
\midrule
Carry.Basic & Fractional-part maps, one-step carries, and the cocycle identity.\\
Carry.Rectangle & Potentials, edge labels, plaquette compatibility, and path
independence.\\
Carry.Cylinders & Integer quotient formula and the exact $Q$-pattern bijection.\\
Carry.Global & Nested intervals and inverse-limit reconstruction.\\
Carry.Furstenberg & Interface theorem stating tail density of the $\times2,\times3$
orbit.\\
Counterexample.Smooth & Formulas for $\nu_{r,s}$, rational coordinates, reduced
denominator exponents, and complete residues.\\
Counterexample.Gaps & Residual bounds, digit gaps, and hitting times.\\
Digits.Lift & Definitions of $L,R$, power inequalities, composition, and valuations.\\
Digits.MulCarry & Signed multiplicativity and the carry-gain expansions.\\
Counterexample.Defects & $D_n,E_n$, sandwiches, valuations, and carry recurrences.\\
Measure.SmoothBC & Almost-everywhere vanishing of
$\omega_{2,3}$.\\
\bottomrule
\end{tabularx}
\end{center}

The dependency graph is
\[
 \texttt{Basic}\to\texttt{Rectangle}\to\texttt{Cylinders}\to
 \begin{cases}
  \texttt{Global},\\
  \texttt{SmoothLift}\leftarrow\texttt{Furstenberg},
 \end{cases}
\]
with the digit modules forming a second branch that meets
\texttt{LeadingGaps} in \texttt{Defects}.

\subsection{Finite theorem signatures}

The following signatures are schematic Lean 4, intended to fix interfaces rather
than claim compilation in the present environment.

\begin{lstlisting}[caption={Schematic finite carry API}]
def potential (R S a b k : Nat) : Nat :=
  k / (2^(R-a) * 3^(S-b))

def hCarry (R S a b k : Nat) : Nat :=
  potential R S (a+1) b k - 2 * potential R S a b k

def vCarry (R S a b k : Nat) : Nat :=
  potential R S a (b+1) k - 3 * potential R S a b k

structure CarryRectangle (R S : Nat) where
  h : Fin R -> Fin (S+1) -> Fin 2
  v : Fin (R+1) -> Fin S -> Fin 3
  plaquette : forall a b,
    3 * (h a b : Nat) + (v a.succ b : Nat) =
    2 * (v a b : Nat) + (h a b.succ : Nat)

theorem cylinder_equiv_carryRectangle (R S : Nat) :
  Fin (2^R * 3^S) <-> CarryRectangle R S
\end{lstlisting}

A convenient proof strategy is to define the forward map by quotient digits and the
inverse map by the upper-right potential.  Path independence can be proved either by
induction on rectangles or by showing that adjacent path swaps preserve the result.
The quotient formula then avoids most real-floor reasoning.

\subsection{Schematic digit API}

\begin{lstlisting}[caption={Schematic digit-lift statements}]
def lift23 (m : Nat) : Nat :=
  (Nat.digits 2 m).eval 3

def lift32 (n : Nat) : Nat :=
  (Nat.digits 3 n).eval 2

theorem lift23_le_rpow (m : Nat) :
  (lift23 m : Real) <= (m : Real) ^ (Real.log 3 / Real.log 2)

theorem lift32_ge_rpow (n : Nat) :
  (n : Real) ^ (Real.log 2 / Real.log 3) <= (lift32 n : Real)

theorem lift23_mul_super (a b : Nat) :
  lift23 a * lift23 b <= lift23 (a*b)

theorem lift32_mul_sub (a b : Nat) :
  lift32 (a*b) <= lift32 a * lift32 b

theorem lift32_lift23 (m : Nat) :
  lift32 (lift23 m) = m
\end{lstlisting}

The formal proof should separate the combinatorial carry-normalization lemma from the
real-power inequalities.  That makes the signed multiplicativity theorem reusable for
general pairs of bases.

\subsection{External interfaces and present library boundary}

As of 22 August 2026, public mathlib documentation contains substantial APIs for
floors, fractional parts, measure theory, and Borel--Cantelli, and it contains density
results for additive irrational rotations on the circle
\cite{MathlibFloor2026,MathlibBorelCantelli2026,MathlibDenseCircle2026}.  The search
used for this report did not locate a public formalization of Furstenberg's
multiplicative $\times2,\times3$ density theorem.  Nor is the full Baker--W\"ustholz
theorem currently a routine imported dependency; a public Lean evaluation problem
still presents it with a placeholder proof \cite{LeanBaker2026}.

Accordingly, the recommended formalization discipline is:
\begin{enumerate}
\item formalize the finite carry bijection and digit-lift package with no axioms;
\item state Furstenberg density behind a single clearly marked interface theorem;
\item state the needed polynomial logarithmic-form bound behind a second interface;
\item derive every counterexample consequence from those two interfaces;
\item replace the interfaces only when independently reviewed formalizations become
available.
\end{enumerate}
This yields a useful machine-checked finite core without pretending that the deep
external theorems have already been imported.

\subsection{Suggested pull-request order}

\begin{enumerate}
\item \texttt{Carry/Cylinders.lean}: quotient digits, plaquettes, exact cardinality.
\item \texttt{Digits/Lift.lean}: order, power bounds, compositions, valuations.
\item \texttt{Digits/MulCarry.lean}: normalization carries and signed products.
\item \texttt{Carry/Realize.lean}: floor realization on half-open intervals.
\item \texttt{Counterexample/FiniteConsequences.lean}: block and residue formulas,
parameterized by an abstract density hypothesis.
\item \texttt{Measure/SmoothBC.lean}: the metric theorem, using mathlib's
Borel--Cantelli infrastructure.
\end{enumerate}
The first three items are finite or elementary and can be reviewed independently of
the open conjecture.

\section{Prioritized research program}
\label{sec:research-program}

The new analysis narrows the next work to a small number of sharply testable
statements.  They are ranked by potential mathematical leverage rather than ease.

\subsection{Priority 1: denominator-sensitive acceleration}

The decisive target remains \Cref{target:subpolynomial-hitting}.  A more structured
version is the following.

\begin{problem}[Paired rational-orbit acceleration]
Let $M,N\ge2$ be integers and suppose
$\log M/\log2=\log N/\log3=\beta\notin\Q$.  Prove that the rational orbit points
\[
 \left(\frac{M^q}{2^{\floor{q\beta}}},
       \frac{N^q}{3^{\floor{q\beta}}}\right),
 \qquad q=2^r3^s,
\]
enter the box
$[1,2^{1/Q})\times[1,3^{1/Q})$ with $q=Q^{o(1)}$ along an unbounded smooth sequence
$Q$.
\end{problem}

This formulation exposes the exact arithmetic unavailable for a generic point: both
coordinates are rational, their numerator exponents are the same smooth $q$, and
their denominator exponents are equal.

\subsection{Priority 2: a joint leading-gap theorem}

\begin{problem}[Matched leading-gap rigidity]
Rule out, for multiplicatively independent output pairs $(M,2)$ and $(N,3)$ linked
by the same $\beta$, simultaneous representations
\[
 M^q=2^\nu+\Delta_2,
 \qquad
 N^q=3^\nu+\Delta_3
\]
with $q=2^r3^s$, the same $\nu$, and
\[
 \Delta_2<2^\nu/Q,
 \qquad
 \Delta_3<3^\nu/Q
\]
for a recurrence scale $q$ below a fixed power of $Q$.
\end{problem}

Single-base digit theorems are insufficient.  The proof would have to use the common
leading position, perhaps through a mixed resultant, a common prime divisor of
suitably normalized residuals, or a subspace theorem with a genuinely coupled
height.

\subsection{Priority 3: high-position multiplication carries}

The carry recurrences suggest a finite combinatorial target.  Let
$\deg_3 C_2(a,b)$ denote the highest position carrying positive weight in
\eqref{eq:binary-carry-gain}, and define the analogous quantity for $C_3$.

\begin{problem}[Paired high-carry lower bound]
Use the fact that $a=M^q$ and $b=N^q$ are synchronized powers to prove that at least
one of the multiplication normalizations for $(M^q)^m$ or $(N^q)^m$ incurs a carry
within $o(\log Q)$ positions of the leading digit, uniformly for some bounded
$m\ge2$, whenever $q$ begins an all-zero rectangle of modulus $Q$.
\end{problem}

A carry that high would give a lower bound conflicting with
\Cref{cor:carry-sparsity}.  The unpaired statement is false because a number may
start with an arbitrarily long zero gap.

\subsection{Priority 4: exploit selected smooth denominator residues}

\Cref{thm:denominator-residues} lets one choose the common top correction
$k\bmod Q$ arbitrarily.  The next useful theorem would attach a height or valuation
cost to special choices, for example $k=0$, $1$, or $Q-1$.  Candidate invariants are:
\begin{itemize}
\item the Kummer class of the reduced rational point modulo $Q$th powers;
\item the Smith profile of the numerator--denominator exponent pair;
\item primitive prime divisors of the paired residuals;
\item the least nonzero multiplication-carry position of the digit lifts.
\end{itemize}
A residue theorem without one of these costs is already exhausted by density.

\subsection{Priority 5: sharpen effective dynamics only where it crosses Baker}

A general improvement from triple exponential to double exponential would be
substantial dynamics, but it would not close the present argument.  The strategically
relevant milestone is a polynomial upper bound
$H_\beta(Q)\ll Q^C$ with an exponent $C$ that can be compared to the explicit Baker
lower exponent.  The decisive milestone is subpolynomial.  Computational and
analytic effort should report results in these terms rather than only in terms of a
smaller density-radius constant.

\subsection{Priority 6: formalize the finite core now}

The carry bijection and digit-lift package are self-contained, nontrivial, and useful
regardless of the conjecture.  Formalizing them now would prevent indexing errors,
make the no-go theorem reusable, and give a stable target against which any claimed
breakthrough can be tested.  In particular, a rival Lean development should be asked
which theorem bridges finite compatibility to arithmetic acceleration; the module
graph in \Cref{sec:lean-blueprint} makes that question precise.

\section{General two-base carry rectangles}
\label{sec:general-bases}

The finite combinatorics is not special to $2$ and $3$.  Recording the general form
clarifies which part of the approach is universal and which part uses the arithmetic
of the conjecture.

Let $a,b\ge2$ be integers.  For $t\in[0,1)$ define
\[
 P_{i,j}(t)=\floor{a^ib^jt},
\]
with horizontal digits
\[
 h_{i,j}=P_{i+1,j}-aP_{i,j}\in\{0,\ldots,a-1\}
\]
and vertical digits
\[
 v_{i,j}=P_{i,j+1}-bP_{i,j}\in\{0,\ldots,b-1\}.
\]

\begin{theorem}[General rectangular carry theorem]
\label{thm:general-carry}
For an $R\times S$ rectangle the edge digits satisfy
\begin{equation}\label{eq:general-plaquette}
 b h_{i,j}+v_{i+1,j}=a v_{i,j}+h_{i,j+1}.
\end{equation}
Conversely, every edge labeling in the indicated digit ranges satisfying all
plaquette equations is realized on a unique cylinder
\[
 \left[\frac{k}{a^Rb^S},\frac{k+1}{a^Rb^S}\right),
 \qquad 0\le k<a^Rb^S.
\]
Consequently there are exactly $a^Rb^S$ compatible rectangles, and their fraction
among all edge labelings is
\begin{equation}\label{eq:general-compatibility-fraction}
 (ab)^{-RS}.
\end{equation}
\end{theorem}

\begin{proof}
The two paths around a plaquette give
\[
 P_{i+1,j+1}=abP_{i,j}+b h_{i,j}+v_{i+1,j}
\]
after moving horizontally first, and
\[
 P_{i+1,j+1}=abP_{i,j}+a v_{i,j}+h_{i,j+1}
\]
after moving vertically first.  This proves \eqref{eq:general-plaquette}.  The
converse repeats the path-independence proof of \Cref{thm:carry-bijection}.  On the
cylinder indexed by $k$,
\[
 P_{i,j}=\floor{\frac{k}{a^{R-i}b^{S-j}}}.
\]
The compatible count is $a^Rb^S$, whereas the unconstrained count is
$a^{R(S+1)}b^{(R+1)S}$; their quotient is $(ab)^{-RS}$.
\end{proof}

For general multiplicatively independent $a,b$, Furstenberg's theorem again makes
the finite compatible language universal.  What is special in Alaoglu--Erd\H{o}s is
that the same exponent produces integral outputs at the two original bases and hence
the paired residual and digit-lift structures of
\Cref{sec:leading-gaps,sec:digit-lifts}.

\section{Conclusion}
\label{sec:conclusion}

This continuation develops a new multiplicative-orbit framework without repeating
the determinant, interpolation, $p$-adic, and additive-word machinery of the supplied
report.  The exact progress is:

\begin{enumerate}
\item a complete classification of every finite $\times2,\times3$ floor-carry
rectangle and its exact entropy collapse;
\item a Furstenberg-based theorem that every compatible finite pattern, and every
smooth floor residue, occurs for every irrational exponent;
\item the resulting no-go theorem for all bounded-window carry strategies;
\item synchronized complete residue classes for the two reduced denominator
exponents of a hypothetical counterexample;
\item arbitrarily long paired leading-zero gaps in powers of the two integer outputs;
\item a precise polynomial-to-triple-exponential corridor for the first such block;
\item a clean subpolynomial hitting-time lemma which would close the conjecture;
\item a digit-lift theory with strict power inequalities, signed multiplicativity,
valuation transfer, positive defect sequences, and exact defect--carry recurrences;
\item a finite Lean architecture separating elementary combinatorics from the two
deep external inputs.
\end{enumerate}

The strongest negative conclusion is as important as the positive ones: finite
symbolic compatibility is completely universal.  The proof cannot be hidden in a
local carry pattern, because every irrational exponent has the same finite carry
language.  The strongest positive conclusion is the location of the missing bridge:
simultaneous integrality must force an arithmetic acceleration of a topologically
generic orbit, or an equivalent high-position carry/leading-gap contradiction.

No full proof is claimed.  The report leaves one central quantitative target rather
than a diffuse collection of analogies.  Any purported breakthrough can now be
stress-tested by asking where it proves one of the following genuinely new facts:
subpolynomial smooth recurrence, matched leading-gap rigidity, or a paired
high-position carry lower bound.  Without such a bridge, the argument remains on the
universal side of the Furstenberg boundary.

\begin{thebibliography}{99}

\bibitem{UnifiedReport2026}
\emph{Alaoglu--Erd\H{o}s unified working report},
attached LaTeX source supplied for this continuation, 20,483 lines; the
SHA-256 checksum is recorded in the source comments,
accessed 22 August 2026.
% Source checksum: eae4567d1f58c0a73594e0b26f93088d6370d7086fef9dbc57414a15919e7877

\bibitem{AE1944}
L.~Alaoglu and P.~Erd\H{o}s,
\newblock \emph{On highly composite and similar numbers},
\newblock Transactions of the American Mathematical Society \textbf{56} (1944),
448--469.
\newblock DOI:
\href{https://doi.org/10.1090/S0002-9947-1944-0011087-2}
{10.1090/S0002-9947-1944-0011087-2}.

\bibitem{Furstenberg1967}
H.~Furstenberg,
\newblock \emph{Disjointness in ergodic theory, minimal sets, and a problem in
Diophantine approximation},
\newblock Mathematical Systems Theory \textbf{1} (1967), 1--49.
\newblock DOI:
\href{https://doi.org/10.1007/BF01692494}{10.1007/BF01692494}.

\bibitem{Boshernitzan1994}
M.~D.~Boshernitzan,
\newblock \emph{Elementary proof of Furstenberg's Diophantine result},
\newblock Proceedings of the American Mathematical Society \textbf{122} (1994),
67--70.
\newblock DOI:
\href{https://doi.org/10.1090/S0002-9939-1994-1195714-X}
{10.1090/S0002-9939-1994-1195714-X}.

\bibitem{BLMV2009}
J.~Bourgain, E.~Lindenstrauss, P.~Michel, and A.~Venkatesh,
\newblock \emph{Some effective results for $\times a\times b$},
\newblock Ergodic Theory and Dynamical Systems \textbf{29} (2009), 1705--1722.
\newblock DOI:
\href{https://doi.org/10.1017/S0143385708000898}
{10.1017/S0143385708000898}.

\bibitem{GayfulinMoshchevitin2023}
D.~Gayfulin and N.~Moshchevitin,
\newblock \emph{On Furstenberg's Diophantine result},
\newblock Moscow Journal of Combinatorics and Number Theory \textbf{12} (2023),
259--272.
\newblock DOI:
\href{https://doi.org/10.2140/moscow.2023.12.259}
{10.2140/moscow.2023.12.259};
\href{https://arxiv.org/abs/2301.08212}{arXiv:2301.08212}.

\bibitem{BakerWustholz1993}
A.~Baker and G.~W\"ustholz,
\newblock \emph{Logarithmic forms and group varieties},
\newblock Journal f\"ur die reine und angewandte Mathematik \textbf{442} (1993),
19--62.
\newblock DOI:
\href{https://doi.org/10.1515/crll.1993.442.19}
{10.1515/crll.1993.442.19}.

\bibitem{AnthropicFable2026}
Anthropic,
\newblock \emph{Claude Fable 5 and Claude Mythos 5},
\newblock announcement dated 9 June 2026, with availability update dated
1 July 2026.
\newblock \url{https://www.anthropic.com/news/claude-fable-5-mythos-5}.

\bibitem{Gao2026}
S.~Gao,
\newblock \emph{Counterexamples to the Jacobian conjecture in dimensions greater
than two},
\newblock arXiv:2608.00222 [math.AG], submitted 31 July 2026.
\newblock \url{https://arxiv.org/abs/2608.00222}.

\bibitem{OEISHadamard2026}
OEIS Foundation,
\newblock \emph{A007299: orders for which a Hadamard matrix is known},
historical note updated in 2026.
\newblock \url{https://oeis.org/A007299}.

\bibitem{Rank30Signal2026}
Yestino--The Signal,
\newblock \emph{An elliptic curve of rank $\ge30$},
public signal page consulted 22 August 2026; the page itself labels the item as not
yet corroborated.
\newblock \url{https://yestino.com/events/an-elliptic-curve-of-rank-30-5eebdc}.

\bibitem{TechCrunchLean2026}
TechCrunch,
\newblock \emph{An unreleased Anthropic model made progress on one of math's biggest
unsolved problems},
\newblock 11 August 2026.
\newblock
\href{https://techcrunch.com/2026/08/11/an-unreleased-anthropic-model-made-progress-on-one-of-maths-biggest-unsolved-problems/}{TechCrunch article page}.

\bibitem{MathlibFloor2026}
mathlib contributors,
\newblock \emph{Mathlib.Algebra.Order.Floor.Defs},
public mathlib4 documentation consulted 22 August 2026.
\newblock
\url{https://leanprover-community.github.io/mathlib4_docs/Mathlib/Algebra/Order/Floor/Defs.html}.

\bibitem{MathlibBorelCantelli2026}
mathlib contributors,
\newblock \emph{Mathlib.Probability.Martingale.BorelCantelli},
public mathlib4 documentation consulted 22 August 2026.
\newblock
\url{https://leanprover-community.github.io/mathlib4_docs/Mathlib/Probability/Martingale/BorelCantelli}.

\bibitem{MathlibDenseCircle2026}
mathlib contributors,
\newblock \emph{Mathlib.Topology.Instances.AddCircle.DenseSubgroup},
public mathlib4 documentation consulted 22 August 2026.
\newblock
\url{https://leanprover-community.github.io/mathlib4_docs/Mathlib/Topology/Instances/AddCircle/DenseSubgroup.html}.

\bibitem{LeanBaker2026}
Lean Formalization Evaluation,
\newblock \emph{Baker--W\"ustholz theorem on linear forms in logarithms},
public problem page consulted 22 August 2026; displayed declaration contains a
placeholder proof.
\newblock
\url{https://lean-lang.org/eval/problems/bakerWustholz_linearForms_logs/}.

\end{thebibliography}

\end{document}
